% ==================///==================///==================///
% ==================/// PREAMBOLO
% ==================///==================///==================///
% I metadati del PDF (pdfauthor, pdftitle) vanno cambiati qui, nella \documentclass.
% Gli stessi nomi compaiono anche in \author, più sotto: tienili allineati.
\documentclass[hyperref={pdfauthor={CHANGE: author 1, CHANGE: author 2, CHANGE: author 3}, pdftitle={CHANGE: title}}, aspectratio={169}]{beamer}
\usepackage[english]{babel}

\usepackage{xcolor}
\usepackage{graphicx}
\usepackage{pifont}
\usepackage{array}
\usepackage{multirow}
\usepackage{adjustbox}
\usepackage{anyfontsize}
\usepackage{setspace}
\usepackage{changepage}
\usepackage{soul, color}
\usepackage[normalem]{ulem}
\usepackage[shortcuts]{extdash}
\usepackage[outline]{contour}
\usepackage{amsmath,amssymb,amsfonts}
\usepackage{mathtools}
\usepackage{xifthen}
\usepackage{csquotes}
\usepackage{makecell}
% \usepackage[maxnames=9, firstinits=true]{biblatex}
% \addbibresource{bibliography.bib}

\usepackage{colortbl}
\usepackage{pbox}

\usepackage{booktabs}

\usepackage{tikz}
\usetikzlibrary{arrows.meta, backgrounds, calc, fit, patterns, positioning, quotes, shapes, shapes.geometric, shapes.misc, shapes.symbols, decorations.pathreplacing}
\usepackage{pgfplots}
\pgfplotsset{compat=newest}
\usepgfplotslibrary{groupplots}
\usepackage{pgfplotstable}
\usepackage{tikzpeople}

\usetheme{Nicola}
% Le slide dell'appendice (Examples) non contano nel totale della talk (es. "7 / 11").
\usepackage{appendixnumberbeamer}

% ==================///==================///==================///
% ==================/// METADATI DELLA PRESENTAZIONE
% ==================///==================///==================///
\title{CHANGE: title}

% \author[forma corta]{\texorpdfstring{testo visibile}{testo semplice per i segnalibri del PDF}}
% - nel testo visibile il relatore è in grassetto e ogni autore ha il numero \inst{N} del suo istituto
% - sotto i nomi, in piccolo, email e pagina web del relatore
\author[CHANGE: author 1, CHANGE: author 2, CHANGE: author 3]{%
    \texorpdfstring{
        \textbf{CHANGE: author 1}\inst{1}, CHANGE: author 2\inst{1}, CHANGE: author 3\inst{2}\\[0.5em]
        \scriptsize \href{mailto:CHANGE:author-1-email@example.org}{CHANGE: author 1 email}\\ \href{https://CHANGE:author-1-web-page}{CHANGE: author 1 web page} }{%
        CHANGE: author 1, CHANGE: author 2, CHANGE: author 3 }%
}

% I numeri \inst{N} devono coincidere con quelli usati in \author.
\institute{%
    \inst{1} CHANGE: affiliation 1 \and
    \inst{2} CHANGE: affiliation 2
}

% Qui la data è il nome dell'evento e l'anno.
\date{CHANGE: event name and year}


\begin{document}

% ==================///==================///==================///
% ==================/// FRAME DEL TITOLO
% ==================///==================///==================///
% Sotto \titlepage c'è un \vfill vuoto: nel deck ICWS di Nicola qui vanno i loghi dei progetti
% e, in piccolo, i ringraziamenti ai finanziamenti.
\begin{frame}
    \titlepage

    \vfill

\end{frame}

% ==================///==================///==================///
% ==================/// MACRO E STILI TIKZ
% ==================///==================///==================///
% Copiati da Nicola (SecAssure 2026): non modificare. Le figure in assets/tex/ ne dipendono.
\newcommand*{\prop}{\ensuremath{p}}
\newcommand*{\propertyNameStyle}[1]{\textsl{#1}}
\newcommand*{\target}{\ensuremath{ToC}}
\newcommand*{\setStyle}[1]{\ensuremath{\mathcal{#1}}}
\newcommand*{\targetSet}{\referred{\setStyle{TOC}}}
\newcommand*{\mechanism}{\ensuremath{m}}
% \newcommand*{\evidCollectionModel}{\ensuremath{\mathcal{E}}}
\newcommand*{\evidCollectionModel}{\ensuremath{E}}
\newcommand*{\evalFn}{\ensuremath{\mathcal{F}}}
% \newcommand*{\certModel}{\ensuremath{\mathcal{CM}}}
\newcommand*{\certModel}{\ensuremath{CM}}
\newcommand*{\certModelSet}{\referred{\setStyle{\certModel}}}
%\newcommand*{\cert}{\ensuremath{\mathcal{C}}}
\newcommand*{\cert}{\ensuremath{C}}
\newcommand*{\certSet}{\referred{\setStyle{\cert}}}
\newcommand*{\evid}{\ensuremath{ev}}
%\newcommand*{}{def}
\newcommand*{\attr}{\ensuremath{attr}}
\newcommand*{\referred}[1]{\ensuremath{\overline{#1}}}

\newcommand*{\repo}{\ensuremath{\mathcal{R}}}

\newcommand*{\signature}{\texttt{sig}}
\newcommand*{\signatureCA}[1][]{\signature\ensuremath{_{\text{CA}_{\text{#1}}}}}
\newcommand*{\signatureAL}[1][]{\signature\ensuremath{_{\text{AL}_{\text{#1}}}}}

\newcommand*{\dimension}{\ensuremath{d}}
\newcommand*{\angled}[1]{\ensuremath{\langle #1\rangle}}


\newcommand{\verificationStep}{\ensuremath{v}}

\newcommand{\targetInput}{\ensuremath{\mathcal{I}}}

%\newcommand{\dimension}{\ensuremath{d}}

\newcommand{\abbrvPreTraining}{\text{\textit{pt}}}
\newcommand{\abbrvAdaptation}{\text{\textit{ad}}}
\newcommand{\abbrvInference}{\text{\textit{in}}}
\newcommand{\abbrvArtifact}{\text{\textit{ar}}}

\newcommand{\viewPreTraining}{\ensuremath{\view_{\abbrvPreTraining}}}
\newcommand{\viewAdaptation}{\ensuremath{\view_{\abbrvAdaptation}}}
\newcommand{\viewInference}{\ensuremath{\view_{\abbrvInference}}}
\newcommand{\viewArtifact}{\ensuremath{\view_{\abbrvArtifact}}}
\newcommand{\viewAdaptationInference}{\ensuremath{\view_{\abbrvAdaptation+\abbrvInference}}}
\newcommand{\dimensionPreTraining}{\ensuremath{\dimension_{\abbrvPreTraining}}}
\newcommand{\dimensionAdaptation}{\ensuremath{\dimension_{\abbrvAdaptation}}}
\newcommand{\dimensionInference}{\ensuremath{\dimension_{\abbrvInference}}}
\newcommand{\dimensionAdaptationInference}{\ensuremath{\dimension_{\abbrvAdaptation+\abbrvInference}}}
\newcommand{\dimensionArtifact}{\ensuremath{\dimension_{\abbrvArtifact}}}

\newcommand{\succSymbol}{\ensuremath{\text{\textit{Succ}}}}
\newcommand{\succFn}[1]{\ensuremath{\succSymbol(#1)}}

\newcommand{\foundation}{\text{found}}

\newcommand{\tableVGap}{.5em}

\newcommand{\steponly}[1]{(#1)}
\newcommand{\step}[1]{step \steponly{#1}}
\newcommand{\steps}[2]{steps \steponly{#1}--\steponly{#2}}
\newcommand{\Step}[1]{Step \steponly{#1}}
\newcommand{\Steps}[2]{Steps \steponly{#1}--\steponly{#2}}

\newcommand{\stepExperimentOnly}[1]{\emph{#1}}
%\newcommand{\stepExperiment}[1]{\emph{step}~\stepExperimentOnly{#1}}
\newcommand{\StepExperiment}[1]{\emph{Step}~\stepExperimentOnly{#1}}

\newcommand*{\gapStyle}[1]{#1}
\newcommand*{\gap}[1]{G#1}

\newcommand{\testCaseOutput}{\ensuremath{out}}
\newcommand{\testCaseResult}{\ensuremath{res}}


%\newcommand{\hypergraphEvalRule}{\ensuremath{H_{\text{eval}}}}\
\newcommand{\hypergraphEvalRule}{\ensuremath{\mathcal{R}}}
\newcommand{\hypergraphAll}{\ensuremath{H_{\text{cert}}}}
\newcommand{\hypergraphInstance}{\ensuremath{\mathcal{H}}}
\newcommand{\verificationPath}{\ensuremath{\mathcal{V}}}

\newcommand{\hypergraphCertModel}{\ensuremath{H}}
%\newcommand{\hypergraphCert}{\ensuremath{H_{\text{cert}}}}
\NewDocumentCommand{\hypergraphCert}{o}{%
  \ensuremath{H_{\text{cert}\IfValueT{#1}{, #1}}}%
}

\newcommand{\hyperpath}{\ensuremath{\text{\emph{path}}}}

\newcommand{\FPathStarSymbol}{\ensuremath{\mathcal{F}}}
\newcommand{\BPathStarSymbol}{\ensuremath{\mathcal{B}}}
\newcommand{\BFPathStarSymbol}{\ensuremath{\mathcal{BF}}}
\newcommand{\FPathStarOf}[1]{\ensuremath{\FPathStarSymbol{(#1)}}}
\newcommand{\BPathStarOf}[1]{\ensuremath{\BPathStarSymbol{(#1)}}}
\newcommand{\BFPathStarOf}[1]{\ensuremath{\BFPathStarSymbol{(#1)}}}
\newcommand{\PathStarSymbol}{\ensuremath{\text{\emph{PathsSimple}}}}
\newcommand{\PathStarOf}[1]{\ensuremath{\PathStarSymbol{(#1)}}}

\newcommand{\evalNode}{\ensuremath{\evalFn}}
\newcommand{\evalNodeAll}{\ensuremath{\widehat{\evalFn}}}

% \newcommand{\nodes}{\ensuremath{V}}
% \newcommand{\edges}{\ensuremath{E}}
\newcommand{\nodes}{\ensuremath{N}}
\newcommand{\arcs}{\ensuremath{A}}

%\newcommand{\node}{\ensuremath{n}}

\newcommand{\modelName}[1]{\texttt{#1}}

\newcommand{\llmGemma}[1][short]{%
  \ifthenelse{\equal{#1}{short}}{%
    \modelName{gemma3:27b}%
  }{%
    \modelName{gemma3:27b-it-q4\_K\_M}%
  }%
}
\newcommand{\llmLlama}[1][short]{%
  \ifthenelse{\equal{#1}{short}}{%
    \modelName{llama3.3:70b}%
  }{%
    \modelName{llama3.3:70b-instruct-q5\_K\_M}%
  }%
}
\newcommand{\llmDeepseek}[1][short]{%
  \ifthenelse{\equal{#1}{short}}{%
    \modelName{deepseek-r1:70b}%
  }{%
    \modelName{deepseek-r1:70b}%
  }%
}
\newcommand{\llmMistral}[1][short]{%
  \ifthenelse{\equal{#1}{short}}{%
    \modelName{mistral:123b}%
  }{%
    \modelName{mistral-large:123b-instruct-2411-q5\_K\_M}%
  }%
}
\newcommand{\llmQwenCoder}[1][short]{%
  \ifthenelse{\equal{#1}{short}}{%
    \modelName{qwen3-coder:30b}%
  }{%
    \modelName{qwen3-coder:30b-a3b-q4\_K\_M}%
  }%
}
\newcommand{\llmGptOss}[1][short]{%
  \ifthenelse{\equal{#1}{short}}{%
    \modelName{gpt-oss:20b}%
  }{%
    \modelName{gpt-oss:20b}%
  }%
}
\newcommand{\llmCodeLlama}[1][short]{%
  \ifthenelse{\equal{#1}{short}}{%
    \modelName{codellama:70b}%
  }{%
    \modelName{codellama:70b-instruct-q5\_K\_M}%
  }%
}

% \newcommand{\llmGemma}{\modelName{gemma3:27b-it-q4\_K\_M}}
% \newcommand{\llmLlama}{\modelName{llama3.3:70b-instruct-q5\_K\_M}}
% \newcommand{\llmDeepseek}{\modelName{deepseek-r1:70b}}
% \newcommand{\llmMistral}{\modelName{mistral-large:123b-instruct-2411-q5\_K\_M}}
% \newcommand{\llmQwenCoder}{\modelName{qwen3-coder:30b-a3b-q4\_K\_M}}
% \newcommand{\llmGptOss}{\modelName{gpt-oss:20b}}
% \newcommand{\llmCodeLlama}{\modelName{codellama:70b-instruct-q5\_K\_M}}


\newcommand{\E}{\ensuremath{E}}
\newcommand{\Efirst}{\ensuremath{\E_{\text{first}}}}
\newcommand{\Elast}{\ensuremath{\E_{\text{last}}}}
\newcommand{\Erandom}{\ensuremath{\E_{\text{random}}}}
\newcommand{\Enull}{\ensuremath{\E_{\text{null}}}}

%\newcommand{\Prompts}{\ensuremath{\mathbb{P}}}

\newcommand{\llm}{\ensuremath{\text{\emph{app}}}}

\newcommand{\probe}{\ensuremath{pr}}
\newcommand{\myN}{\ensuremath{n}}
\newcommand{\myM}{\ensuremath{m}}
\newcommand{\accordance}{\ensuremath{\text{\emph{res}}}}
%\newcommand{\probeSet}{\ensuremath{P}}
\newcommand{\prompt}{\ensuremath{\text{\emph{prompt}}}}

\newcommand{\answer}{\ensuremath{\text{\emph{ans}}}}

\newcommand{\orderBias}{\ensuremath{\text{\emph{bias}}}}

\newlength{\oldtabcolsep}
\setlength{\oldtabcolsep}{\tabcolsep}

\newcommand{\tailSetSymbol}{\ensuremath{\text{\emph{tail}}}}
\newcommand{\headSetSymbol}{\ensuremath{\text{\emph{head}}}}

\newcommand{\node}{\ensuremath{\text{\emph{node}}}}
\newcommand{\arc}{\ensuremath{\text{\emph{arc}}}}

\newcommand{\headSet}[1]{\ensuremath{\headSetSymbol(#1)}}
\newcommand{\tailSet}[1]{\ensuremath{\tailSetSymbol(#1)}}

\newcommand{\guardFunctionSymbol}{\ensuremath{G}}
\newcommand{\guardFunction}[1]{\ensuremath{\guardFunctionSymbol(#1)}}
\newcommand{\evidenceFunctionSymbol}{\ensuremath{EV}}
\newcommand{\evidenceFunction}[1]{\ensuremath{\evidenceFunctionSymbol(#1)}}
\newcommand{\nil}{\texttt{nil}}
\newcommand{\error}{\texttt{error}}

%%%%%%%%%
% new macro

\newcommand{\propertyLabel}{\ensuremath{p}}
\newcommand{\behaviorSet}{\ensuremath{\mathcal{B}}}
\newcommand{\verificationStepTarget}{\ensuremath{t}}
\newcommand{\verificationStepActivity}{\ensuremath{s}}
\newcommand{\conditionGuardSymbol}{\ensuremath{\text{\emph{g}}}}
\newcommand{\conditionGuardFunction}[1]{\ensuremath{\conditionGuardSymbol(#1)}}

%\newcommand{\verificationPathAnnotated}{\ensuremath{\overline{\verificationPath}}}

\newcommand{\verificationGraphAnnotated}{\ensuremath{\overline{\behaviorSet}}}
\newcommand{\verificationNodesAnnotated}{\ensuremath{\overline{\nodes}}}
\newcommand{\verificationArcsAnnotated}{\ensuremath{\overline{\arcs}}}


\newcommand{\sourceFunctionSymbol}{\ensuremath{\text{\emph{source}}}}
\newcommand{\sourceFunction}[1]{\ensuremath{\sourceFunctionSymbol(#1)}}

\newcommand{\firstFunctionSymbol}{\ensuremath{\text{\emph{first}}}}
\newcommand{\firstFunction}[1]{\ensuremath{\firstFunctionSymbol(#1)}}

\newcommand{\lastFunctionSymbol}{\ensuremath{\text{\emph{last}}}}
\newcommand{\lastFunction}[1]{\ensuremath{\lastFunctionSymbol(#1)}}

\newcommand{\arcCounterSymbol}{\ensuremath{c}}
\newcommand{\arcCounterOf}[1]{\ensuremath{\arcCounterSymbol[#1]}}



\newcommand{\applTypeTextStandalone}{Type 1: Standalone}
\newcommand{\applTypeTextIntegratedArtifactGeneration}{Type 2: Integrated -- Artifact Generation}
\newcommand{\applTypeTextIntegratedDecisionMaking}{Type 3: Integrated -- Decision-Making}

\newcommand{\probeOutput}{\ensuremath{\text{\emph{out}}}}
\newcommand{\probeOutputTrue}{\ensuremath{\probeOutput^{\top}}}
\newcommand{\probeOutputFalse}{\ensuremath{\probeOutput^{\bot}}}

\newcommand{\FKGLSymbolFunction}{\ensuremath{\text{\emph{FK}}}}
\newcommand{\FKGLFunction}[1]{\ensuremath{\FKGLSymbolFunction(#1)}}
\newcommand{\DLSymbolFunction}{\ensuremath{\text{\emph{DL}}}}
\newcommand{\DLFunction}[1]{\ensuremath{\DLSymbolFunction(#1)}}

\newcommand{\detail}[1]{\emph{#1}}
\newcommand{\detailLow}{\detail{Low}}
\newcommand{\detailMediumLow}{\detail{Medium--Low}}
\newcommand{\detailMediumHigh}{\detail{Medium-High}}
\newcommand{\detailHigh}{\detail{High}}

% displays a "check mark"
\newcommand{\ok}{\ding{51}}
% displays an "x mark"
\newcommand{\notOk}{\ding{55}}


\newcommand{\bigO}{\mathcal{O}}

\newcommand{\arcsCardinality}{\ensuremath{\vert \arcs \vert}}
\newcommand{\nodesCardinality}{\ensuremath{\vert \nodes \vert}}

\newcommand{\nodesFailed}{\ensuremath{\overline{F}}}


\newcommand{\spaceOne}{~~~~}

\newcommand{\challenge}[1]{\textbf{C#1}}

\newcommand{\parMini}[1]{\parbox{\linewidth}{\tiny #1}}
\newcommand{\fullCiteMini}[1]{\parMini{\fullcite{#1}}}


\newcommand{\parAlmostMini}[1]{\parbox{\linewidth}{\footnotesize #1}}
\newcommand{\fullCiteAlmostMini}[1]{\parAlmostMini{\fullcite{#1}}}

\newcommand{\timestyle}[1]{#1}
\newcommand{\timeshort}{\timestyle{S}}
\newcommand{\timemedium}{\timestyle{M}}
\newcommand{\timelong}{\timestyle{L}}


\newcommand*{\service}[1]{\ensuremath{s_{#1}}}

\newcommand*{\pavail}{\ensuremath{\prop_\text{avail}}}
\newcommand*{\ppriv}{\ensuremath{\prop_{\text{priv}}}}
\newcommand*{\pConfTransit}{\ensuremath{\prop_{\text{conf\_net}}}}
\newcommand*{\pConfStorage}{\ensuremath{\prop_{\text{conf\_str}}}}


%\newcommand*{\relwordTargetCloudTrad}{def}
\newcommand*{\abbrvServiceStorage}{\text{str}}
\newcommand*{\abbrvServiceDatabase}{\text{db}}
\newcommand*{\abbrvServiceAuth}{\text{auth}}
\newcommand*{\abbrvServiceNetwork}{\text{net}}
\newcommand*{\abbrvServiceKube}{\text{k8s}}
\newcommand*{\abbrvServiceAnalytics}{\text{dawp}}
% \newcommand*{\abbrvServiceNetworkInt}{\text{netinf}}
% \newcommand*{\abbrvServiceStorageInt}{\text{strinf}}

\newcommand*{\abbrvServiceNetworkIntCloud}{\text{netinf\_cl}}
\newcommand*{\abbrvServiceStorageIntCloud}{\text{strinf\_cl}}
\newcommand*{\abbrvServiceNetworkIntEdge}{\text{netinf\_edge}}
\newcommand*{\abbrvServiceStorageIntEdge}{\text{strinf\_edge}}

\newcommand*{\abbrvMechanismUserAuth}{\text{user}}
\newcommand*{\abbrvMechanismScheduling}{\text{sched}}
\newcommand*{\abbrvMechanismData}{\text{data}}

\newcommand*{\abbrvHandlerApp}{\text{app}}
\newcommand*{\abbrvHandlerCSP}{\text{csp}}
\newcommand*{\abbrvHandlerISP}{\text{isp}}

\newcommand*{\CA}{\emph{\gls{CA}}}
\newcommand*{\CAApplication}{\CA\ensuremath{_{\abbrvHandlerApp}}}
\newcommand*{\CACSP}{\CA\ensuremath{_{\abbrvHandlerCSP}}}
\newcommand*{\CAISP}{\CA\ensuremath{_{\abbrvHandlerISP}}}

\newcommand*{\AL}{\emph{\gls{AL}}}
\newcommand*{\ALApplication}{\AL\ensuremath{_{\abbrvHandlerApp}}}
\newcommand*{\ALCSP}{\AL\ensuremath{_{\abbrvHandlerCSP}}}
\newcommand*{\ALISP}{\AL\ensuremath{_{\abbrvHandlerISP}}}

\newcommand{\approachSOTA}{\text{\emph{SOTA}}}
\newcommand{\approachOur}{\text{\emph{multi}}}
\newcommand{\certTimeSymbol}{\ensuremath{T}}
\newcommand{\certTimeFunction}[1]{\ensuremath{\certTimeSymbol}(#1)}

\newcommand*{\certStatusStyle}[1]{\texttt{#1}}
\newcommand*{\certStatusValid}[1]{\certStatusStyle{valid(}#1\certStatusStyle{)}} % chktex 9 chktex 36
\newcommand*{\certStatusSuspended}[1]{\certStatusStyle{suspended(}#1\certStatusStyle{)}} % chktex 9 chktex 36
\newcommand*{\certStatusRevoked}[1]{\certStatusStyle{revoked(}#1\certStatusStyle{)}} % chktex 9 chktex 36

\newcommand{\highlightprofilecolor}{LightGoldenRod!100}
\colorlet{HighlightFill}{LightGoldenRod!50}
\colorlet{HighlightFillLight}{LightGoldenRod!50}
\colorlet{HighlightProfile}{LightGoldenRod!100}
\definecolor{NotSoDarkGrey}{RGB}{160, 160, 160}

%%% tikz macro
\colorlet{PoisonedRedForText}{white!70!red!100}
\colorlet{PoisonedRedOut}{red!70}
\colorlet{UnPoisonedGreen}{green!40!black!40}
\colorlet{UnPoisonedGreenForText}{green!30!white!90}
\colorlet{HighlightColor}{Goldenrod}
\colorlet{default}{DarkBlue}
\colorlet{HighlightFill}{LightGoldenRod!60}

\definecolor{Completed}{RGB}{82,163,0}
\definecolor{InProgress}{RGB}{204, 102, 0}

\newcommand{\normalwidth}{0.3pt}
\newcommand{\thickwidth}{0.8pt} %means thick

\newcommand{\forhash}{for hash}
%\newcommand{\forTool}{for tool}
\newcommand{\forrisk}{for risk}

\newcommand{\forToolTraining}{for tool training}
\newcommand{\forToolTesting}{for tool testing}

\newcommand{\poisonedRF}{1}
% determines the actual title of the our approach part of the figures,
% considered only when is \forTool.
% \newcommand{\headerEnsemblePoisoned}{Poisoned Hash\-/Based Ensemble}
% \newcommand{\headerEnsembleNonPoisoned}{Non\-/Poisoned Hash\-/Based Ensemble}
\newcommand{\headerEnsemblePoisoned}{Hash\-/Based Ensemble (Poisoned Training Sets)}
\newcommand{\headerEnsembleNonPoisoned}{Hash\-/Based Ensemble (Clean Training Sets)}
\newcommand{\headerEnsembleDefault}{Hash\-/Based Ensemble}

\newcommand{\headerRfTestingPoisoned}{Poisoned}
\newcommand{\headerRfTestingNonPoisoned}{Non\-/Poisoned}
\newcommand{\headerRfOther}{}

\newcommand{\codePoisoned}{Poisoned}
\newcommand{\codeNonPoisoned}{Non\-/Poisoned}

\newcommand{\devilLarge}{Large}
\newcommand{\devilSmall}{Small}

\tikzset{
  invisible/.style={opacity=0,text opacity=0},
  visible on/.style={alt=#1{}{invisible}},
  alt/.code args={<#1>#2#3}{%
    \alt<#1>{\pgfkeysalso{#2}}{\pgfkeysalso{#3}} % \pgfkeysalso doesn't change the path
  },
  on slide/.code args={<#1>#2}{%
  \only<#1>{\pgfkeysalso{#2}}%
}
}

\tikzset{
    every node/.style={
        font=\scriptsize,
    },
    dataset/.style={
        draw,
        cylinder,
        shape border rotate=90,
        align=center,
    },
    on path/.style={
        above
    },
    rf/.style={
        draw,
        rounded corners,
        %double,
        % NODES MUST BE FILLED UPFRONT, OR
        % ANIMATING FILL IS MEANINGLESS
        fill=white
    },
    test path/.style={
        dashed
    },
    connections/.style={
        very thin
    },
    partition/.style={
        draw,
        cylinder,
        shape border rotate=90,
        align=center,
        minimum height=.6cm,
        minimum width=.35cm,
        fill=white
    },
    link/.style={
        >=Latex
    },
    container/.style={
        very thin,
        draw
    },
    poisoning link/.style={
        link,
        decorate,
        decoration={
            amplitude=.5pt,
            segment length=4pt,
            zigzag},
        ultra thin,
       % >=Stealth
    },
    data point/.style={
        draw,
        %very thin,
        %rectangle,
        rounded corners,
        font={
            \tiny
        },
        fill=white,
        %fill opacity=.5
    },
    data point moving/.style={
        circle,
        draw,
        fill=black,
        inner sep=1pt
    },
    poisoned/.style={
        fill = PoisonedRedOut
    },
    majority voting/.style={
        isosceles triangle,
        draw,
    },
    on slide/.code args={<#1>#2}{%
      \only<#1>{\pgfkeysalso{#2}}%
    },
    highlight base/.style={
        very thick
    },  
    highlight path/.style={
        highlight base,
        draw=LightGoldenRod,
    },
    highlight node/.style={
        highlight base,
        fill=HighlightFill,
    },
    highlight double/.style={
        double=LightGoldenRod,
        double distance=1.2pt,
        line width=.3pt
    } ,
    arrow path unpoisoned/.append style={
        double=UnPoisonedGreen,
        arrows={-Latex[fill=UnPoisonedGreen]}
    },
    arrow path poisoned/.append style={
        double=PoisonedRedOut,
        arrows={-Latex[fill=PoisonedRedOut]}
    },
    % cross line/.style={
    %     %double=scope’s color setting,
    %     double=default,
    %     draw=white,
    %     % the thickness of the outer line,
    %     very thick, 
    %     double distance=.1pt
    % },
    % cross line arrow/.style={
    %     arrows={-Latex[]}
    % }
    highlight path poisoned/.style={},
    highlight path unpoisoned/.style={},
    invisible with opacity/.style={opacity=0.025},
    visible with opacity on/.style={alt=#1{}{invisible with opacity}},
    % alt/.code args={<#1>#2#3}{%
    %   \alt<#1>{\pgfkeysalso{#2}}{\pgfkeysalso{#3}} % \pgfkeysalso doesn't change the path
    % },
    highlight line in pgfplot/.style={
        LightGoldenRod,
        line width=3pt
    },
    poisoning attack/.style={
        draw,
        %dashed
        fill=white
    },
    super container/.style={
        draw,
        dashed
    },
    low link/.style={
        link,
        ultra thin
    },
    bar/.style={
        thin,
        dotted
    },
    from test/.style={
        ultra thin
    },
    highlight accuracy/.style={
        ultra thick
    },
    explanation/.style={
        font={\Large},
    },
    explanation link/.style={
        >=Latex,
    },
    line am/.style={},
    line m2/.style={
        dashed
    }
}

\newcommand{\includegraphicswithwisibility}[3]{
\begin{tikzpicture}
\node[visible with opacity on={#1}] at (0,0) {\includegraphics[#2]{#3}};
\end{tikzpicture}
}

\tikzset{
    entity/.style={
      draw,
      ellipse,
      very thin
    },
    artifact/.style={
      rectangle,
      draw,
      rounded corners,
      fill=VeryLightGrey!50
    },
    service/.style={
      chamfered rectangle,
      chamfered rectangle angle=30,
      draw,
      %fill=VeryLightGrey!120
    },
    every node/.style={
        font={\scriptsize},
        text centered,
    },
    on link/.style={
      fill=white,
    },   
    connector/.style={
        >=stealth,
    },
    connector assurance/.style={
        >=Stealth,
        dashed
    },
    on slide/.code args={<#1>#2}{%
      \only<#1>{\pgfkeysalso{#2}}%
    },
    highlight base/.style={
      very thick
    },
    highlight/.style={
      highlight base,
      fill=LightGoldenRod,
    },
    highlight profile/.style={
      \highlightprofilecolor
    },
    highlight double/.style={
      double=LightGoldenRod
    },
    highlight text/.style={
      text=Goldenrod,
    },
    highlight link/.style={
      %thick,
      color=Goldenrod
    },
    highlight light/.style={
      fill=LightGoldenRod!50
    },
    highlight very light/.style={
      fill=LightGoldenRod!20
    },
      element/.style={
          inner sep=1.5pt,
          font=\footnotesize
      },
      connector/.style={
          >=Stealth,
      },
      lattice layer label/.style={
          text width=2.9cm,
          font={\footnotesize},
      },
      not chosen color/.style={
          black!50
      },
      lattice layer separator/.style={
          dashed,
      },
      service label/.style={
          fill=white,
          circle,
          draw,
          very thick,
          inner sep=2.5pt,
      },
      lattice element bound/.style={
          fill=black!15,
          double
      },
      lattice element selected/.style={
          fill=Goldenrod!15,
          double
      },
      lattice element not chosen/.style={
          not chosen color,
          densely dotted
      },
      lattice element not chosen link/.style={
          not chosen color,
          densely dotted
      },
      service element not chosen/.style={
          text=black!50,
          draw=black!50,
          dashed,
      },
      header/.style={
      },
      main box/.style={
          font={\small}
      },
      link/.style={
        >=Latex,
      },
      activity/.style={
        fill=white
      },
    % for cert history
      arrow/.style={
        draw,
        minimum height=1cm,
        inner sep=2pt,
        shape=signal,
        signal from=west,
        signal to=east,
        signal pointer angle=110,
        top color=DarkBlue,
        bottom color=DarkBlue,
        text=White,
        text width=.24\textwidth,
        font = {\small}
    },
    desc/.style={
      %text width=3.9cm,
      text width=.22\textwidth,
      fill=VeryLightGrey,
      %minimum height=1cm,
      %anchor=south
      % text depth=1cm,
      rectangle,
      font = {\footnotesize}
    },
    highlight line in plot/.style={
      %Goldenrod,
      %line width=3pt
  },
  basic question/.style={
    rectangle callout,
    %fill=HighlightFillLight,
    fill=VeryLightGrey!50,
    draw,
  },
  object/.style={
    draw,
    rectangle,
    rounded corners,
    fill=white
    %minimum height=1.8cm
},
  % this combination is needed because otherwise when we use
  % labels, the underlying drawing is visible through the foreground label
  a bit of transparency/.style={
    %semitransparent,
    opacity=0.5, % it was 0.85
    draw opacity=0.5,
  },
  a bit of transparency label/.style={
    a bit of transparency,
    opacity=0.85,
    text opacity=1
  },
% used as "header" of a box
  box label/.style={
    fill=DarkBlue,
    text=white,
    double,
    %draw=white,
    rounded corners
  },
  basic question no color/.style={
    rectangle callout,
    %basic question no callout pre,
    %fill=HighlightFillLight,
    %fill=white,
    fill=VeryLightGrey!50,
    draw,
    thick
  },
  aggregation block/.style={
    circle,
    draw,
    inner sep=1pt,
    %fill=CellColorOk
  },
    link comment/.style={
      dashed
    },
}
\tikzset{
	entity/.style={
			draw,
			ellipse,
			%circle,
			thin,
			%inner sep=0pt
		},
	artifact/.style={
			rectangle,
			draw,
			rounded corners,
			%fill=VeryLightGrey!50
		},
	service in certification process/.style={
			chamfered rectangle,
			chamfered rectangle angle=30,
			draw,
			%fill=VeryLightGrey!120
		},
	% every node/.style={
	% 		font={\small},
	% 		text centered,
	% 	},
	on link/.style={
			fill=white,
		},
	link/.style={
			>=Latex,
			thin
		},
	probe link/.style={
			link,
			dashdotted
		},
	ca cert/.style={
			entity,
			inner sep=0pt,
		},
	step separator/.style={
			draw,
			thin,
			dotted
		},
	container/.style={
			very thin,
			draw,
			rounded corners,
			%densely dotted
		},
	container line/.style={
			very thin,
			rounded corners,
			%densely dotted
		},
	service/.style={
			rectangle,
			draw,
			fill=white,
			minimum width=.35cm,
			minimum height=.35cm
		},
	label desc/.style={},
	label desc with shape/.style={
			label desc,
			rounded corners,
			draw
		},
	boxes/.style = {
			draw=black,
			rectangle,
		},
	rboxes/.style = {
			boxes,
			rounded corners
		}
}

\tikzset{
	basic graph node/.style={
		circle,
		draw,
		minimum width=.6cm
		%minimum width=\refwidth
	},
	verification step/.style={
		basic graph node,
		fill=white
	},
	evaluation function/.style={
		basic graph node,
		%fill=black
		double,
		fill=white
	},
	evaluation function global/.style={
		draw,
		double,
		rectangle,
		minimum width=.6cm,
		minimum height=.6cm,
		%fill=black
	},
	hyperarc/.style={
		>=Latex,
	},
}

\tikzset{
	pics/hyperarc between fan in/.style n args={5}{
		code={
			\path let \p1 = ($(#1)!.5!(#2)$), \p2 = ($(#2)!.5!(#3)$) in
        		coordinate (#4) at (\x1, \y2);
   			\foreach \src in {#1, #2} {
        		\draw[-, hyperarc, #5] (\src) -- (#4);
    		}
    		\draw[->, hyperarc, #5] (#4) -- (#3);
		}
	},
		pics/hyperarc between fan out/.style n args={5}{
		code={
			\path let \p1 = ($(#1)!.5!(#2)$), \p2 = ($(#2)!.5!(#3)$) in
        		coordinate (#4) at (\x1, \y2);

   			\foreach \src in {#1, #2} {
        		\draw[<-, hyperarc, #5] (\src) -- (#4);
    		}
    		\draw[-, hyperarc, #5] (#4) -- (#3);
		}
	},
}

\tikzset{
	visited node/.style={
		%double
	},
	visited arc/.style={
		thick
	},
	pruned node/.style={
		fill=gray!30
	},
	pruned arc/.style={
		draw=gray!50
	}
}

\tikzset{
	evidence content/.style={
		font=\scriptsize
	}
}

\definecolor{colorBlindBiOrange}{HTML}{E66100}
\definecolor{colorBlindBiViolet}{HTML}{5D3A9B}

\tikzset{
	orange in violet out double/.style={
		fill=colorBlindBiOrange,
		double=colorBlindBiViolet,
		draw=colorBlindBiViolet,
		very thick
	},
	orange out/.style={
		draw=colorBlindBiOrange,
		very thick
	}
}

\tikzset{
	highlighted hyperarc in path/.style={
		thick
	}
}

\tikzset{
	hyperarc in verification path/.style={
		very thick,
        highlight path
	},
	node in verification path/.style={
		very thick,
        highlight path
	},
}

\tikzset{
  speaking question/.style={
    ellipse callout,
    fill=VeryLightGrey!50,
    draw,
    %very thick
    %thick
  },
}

\tikzset{
    certificate callout/.style={
        rectangle callout,
        fill=VeryLightGrey!50,
        draw, 
    }
}

% ==================///==================///==================///
% ==================/// SLIDE
% ==================///==================///==================///
% GUIDA: una decina di slide, al massimo 12, per un intervento di 12 minuti.
% !TEX root = main.tex
% GUIDA: 1-2 slide di contesto (Background e Motivations). Vedi gli esempi
% legacy per le animazioni con figura e scenario.
\begin{frame}
    \frametitle{Context: Big Data APIs \& Security Assurance}

    % Goal della slide: Stabilire il contesto (Big Data API) e far capire subito dove si colloca il problema.
    \begin{itemize}
        \item \textbf{Big Data Ecosystem:} rely on microservices and distributed pipelines, exposed via REST APIs

        \item \textbf{The Trust Boundary:} API Gateways centralize critical enforcement (authentication, authorization, \dots)

        \item \textbf{The Assurance Gap:} Security focuses on infrastructure, while API testing targets functional correctness, leaving the interface layer exposed.

        \item \textbf{The Consequence:} Interface-level vulnerabilities (e.g., cross-tenant access) can compromise the entire Big Data pipeline
    \end{itemize}

\end{frame}

\begin{frame}
    \frametitle{State of Art Limits (Motivations) \& Objectives}

    \begin{columns}[T]
        \begin{column}{.5\textwidth}
            \begin{block}{State of the Art Limitations}
                \begin{itemize}
                    \item \textbf{Syntactic Blindness:} \\
                          Checks request/response syntax, ignoring security contract semantics
                    \item \textbf{The Oracle Problem:} \\
                          Lacks automated criteria for security guarantee.
                    \item \textbf{Unreliable Execution:} \\
                          Non-deterministic outcomes (flakiness) prevent continuous CI/CD integration
                \end{itemize}
            \end{block}
        \end{column}

        \begin{column}{.5\textwidth}
            % Objectives
            \begin{block}{Our Objectives}
                \begin{itemize}
                    \item \textbf{Semantic Verification:} \\
                          Move beyond syntax to verify the actual behavior of security controls.
                    \item \textbf{Specified Oracles:} \\
                          Automate pass/fail verdicts using hardcoded, domain-derived rules.
                    \item \textbf{Agnostic \& Deterministic:} \\
                          Build a reproducible engine, portable across any documented target.
                \end{itemize}
            \end{block}
        \end{column}

    \end{columns}

\end{frame}

\begin{frame}
    \frametitle{CHANGE: Our Approach: Overview}
    % GUIDA: una frase sull'approccio e sotto un punto per componente. Le slide dopo corrispondono ai punti, nello stesso ordine:
    % è un indice implicito dei contenuti, non una scaletta.

    Main \alert{idea}: CHANGE: one sentence stating the idea of the approach

    \vspace{1em}

    \begin{itemize}\setlength{\itemsep}{1em}
        \item CHANGE: aspect 1 of the idea, with the \alert{key words} highlighted
        \item CHANGE: aspect 2 of the idea
        \item CHANGE: aspect 3 of the idea
    \end{itemize}

\end{frame}

\begin{frame}
    \frametitle{CHANGE: Our Approach: Details}
    % GUIDA: una slide per componente, dall'alto verso il basso: dove sta nell'approccio complessivo, poi cos'è. Meglio un'immagine,
    % anche animata, di una definizione lunga. Si copia il frame per ogni componente (Details (1), (2), ...).
    % Per un blocco lungo si può rimpicciolire il testo: \begin{block} ... \end{block} dentro { \footnotesize ... }
    \begin{block}{CHANGE: element name}
        A \alert{CHANGE: element} is a tuple \alert{$\angled{x_1, \{y_1,\ldots, y_n\}, \referred{X}, \signature}$}, where

        \vspace{1em}

        \begin{itemize}\setlength{\itemsep}{.5em}
            \item \alert{$x_1$} is CHANGE: description of element 1
            \item \alert{$\{y_i\}$} is a set of \emph{CHANGE: elements}, CHANGE: description of element 2
            \item \alert{$\referred{X}$$=$$\{\referred{x}_1,$ \ldots$,$ $\referred{x}_n\}$} is a set of \alert{\emph{referred elements}}, \alert{each} $\referred{x}_i$ \alert{CHANGE: relation to the element} (CHANGE: description of element 3)
            \item \alert{$\signature$} is CHANGE: description of element 4.
        \end{itemize}
    \end{block}

\end{frame}

% Nell'approccio può stare anche una figura come quella delle Motivations (01_intro.tex), con i
% colori che evidenziano i passi: stessa struttura (macro dei numeri di overlay definite nel frame,
% figura in assets/tex/, \input dentro adjustbox). Non c'è una slide a parte: si riusa la stessa tecnica.
% GUIDA sull'esempio: si può raccontare dentro ogni componente, mentre lo presenti, oppure passo passo dopo
% l'approccio. Dipende dal tempo.

\begin{frame}
    \frametitle{CHANGE: State diagram}
    % GUIDA: un altro aspetto del processo, mostrato con un diagramma e poche parole (qui un diagramma di stati).

    \begin{adjustbox}{max totalsize={\linewidth}{\textheight},center}
        		\begin{tikzpicture}
			\tikzset{
				every node/.style={font=\scriptsize}
			}

			% \node[rboxes] (valid) {Valid};
			% \node[rboxes] (suspended) [right=2 of valid] {Suspended};
			% \node[rboxes] (revoked) [right=2 of suspended] {Revoked};
			\node[rboxes] (valid) {\certStatusValid{\cert}};
			\node[rboxes] (suspended) [right=2 of valid] {\certStatusSuspended{\cert}};
			\node[rboxes] (revoked) [right=2 of suspended] {\certStatusRevoked{\cert}};

			\draw[-latex] (valid) to [out=30,in=180-30] node [above]{\makecell{AL failure\\\referred{\cert} suspended}} (suspended);
			\draw[-latex] (suspended) to [out=180+30,in=-30] node [below]{\makecell{Issue remediated}} (valid);
			\draw[-latex] (valid.north) -- ++(0,1.5) -- node[midway,above]{\referred{\cert} revoked} ([yshift=1.5cm]revoked.north) -- (revoked.north);
			\draw[-latex] (suspended) -- node [above]{\makecell{Issue not\\remediated}} (revoked);
		\end{tikzpicture}
    \end{adjustbox}

\end{frame}

% !TEX root = main.tex
% GUIDA: slide di esperimenti e risultati: testo con immagine, contenuto
% sostituibile, tabella e grafico.
\begin{frame}
    \frametitle{CHANGE: Experiments}

    CHANGE: obiettivo degli esperimenti e come sono impostati. Se sono tanti, mostra solo il sottoinsieme che fa passare l'intuizione.
    Una o due slide.

\end{frame}

\begin{frame}
    \frametitle{CHANGE: Results}

    CHANGE: i risultati, per esempio come grafici, con una sola conclusione per grafico.

\end{frame}

% !TEX root = main.tex

\begin{frame}
    \frametitle{Conclusions \& Future Work}

    Conclusions
    \begin{itemize}
        \item Contract-driven testing successfully evaluates security semantics beyond simple syntactic conformance (RQ1).
        \item Specified oracles provide deterministic, automated pass/fail criteria.
        \item The engine's reproducibility enables reliable integration into continuous release cycles (RQ2).
    \end{itemize}


    Future Work
    \begin{itemize}
        \item Develop adapters for managed cloud API gateways (AWS, Azure, GCP).
        \item Transition from manually authored rules to automated generation of security oracles.
    \end{itemize}


\end{frame}

\begin{frame}

    % Il blocco "Thanks! Questions?" usa i colori e il carattere del titolo (stesso aspetto della barra
    % del titolo del primo frame). Sotto, i contatti: sono gli stessi di main.tex (nome, email, sito) e si scrivono due volte.
    \begin{center}
        \begin{beamercolorbox}[sep=.75em]{title}
            \centering
            \usebeamerfont{title}Thanks! Questions?%
        \end{beamercolorbox}

        \vspace{2em}

        {\large Enea Manzi}\\
        \vspace{.25em}
        {\small
            \href{mailto:enea.manzi@unimi.it}{enea.manzi@unimi.it}\\
            \href{https://eneamanzi.github.io/}{https://eneamanzi.github.io/}\\
            \href{https://www.linkedin.com/in/eneamanzi}{https://www.linkedin.com/in/eneamanzi}
        }

    \end{center}

\end{frame}


% Esempi: tecniche pronte da copiare nelle slide, in appendice. Per toglierli cancella queste due righe.
\appendix
% !TEX root = main.tex
% ==================///==================///==================///
% ==================/// ESEMPI (appendice, non numerata con la talk)
% ==================///==================///==================///
% Pezzi di codice pronti da copiare nelle slide, uno per tecnica, presi dalle presentazioni di Nicola
% (SecAssure 2026 e ICWS 2024) e resi generici. Ogni file in examples/ ha in cima un commento che
% spiega la tecnica e cosa copiare. Per togliere la sezione basta cancellare, in main.tex, le righe
% che includono questo file (\appendix e % !TEX root = main.tex
% ==================///==================///==================///
% ==================/// ESEMPI (appendice, non numerata con la talk)
% ==================///==================///==================///
% Pezzi di codice pronti da copiare nelle slide, uno per tecnica, presi dalle presentazioni di Nicola
% (SecAssure 2026 e ICWS 2024) e resi generici. Ogni file in examples/ ha in cima un commento che
% spiega la tecnica e cosa copiare. Per togliere la sezione basta cancellare, in main.tex, le righe
% che includono questo file (\appendix e % !TEX root = main.tex
% ==================///==================///==================///
% ==================/// ESEMPI (appendice, non numerata con la talk)
% ==================///==================///==================///
% Pezzi di codice pronti da copiare nelle slide, uno per tecnica, presi dalle presentazioni di Nicola
% (SecAssure 2026 e ICWS 2024) e resi generici. Ogni file in examples/ ha in cima un commento che
% spiega la tecnica e cosa copiare. Per togliere la sezione basta cancellare, in main.tex, le righe
% che includono questo file (\appendix e \input{05_examples.tex}).

\begin{frame}[plain,c]
    \centering
    {\Huge\bfseries \color{DarkBlue}01 Examples}\\[1em]
    {\large Ready-to-copy code for the techniques seen in Nicola's presentations}
\end{frame}

\input{examples/00_basic/02_highlighted_list_next_to_figure.tex}
\input{examples/00_basic/03_technical_definition.tex}
\input{examples/00_basic/04_progressive_content_with_pause.tex}
\input{examples/00_basic/05_replacing_content_with_only.tex}
\input{examples/00_basic/06_tables.tex}
\input{examples/00_basic/07_plots_external_data.tex}
\input{examples/01_advanced/08_synchronized_characters_and_speech.tex}
\input{examples/01_advanced/09_chevrons.tex}
\input{examples/01_advanced/10_reusable_pic_component.tex}
\input{examples/01_advanced/11_callouts_and_speech_bubbles.tex}
\input{examples/01_advanced/12_title_logos_and_funding.tex}

% Le due slide originali di Nicola hanno dipendenze specifiche, caricate
% solo quando queste slide legacy vengono incluse.
\input{examples/legacy-nicola/00_macro_legacy.tex}
\input{examples/legacy-nicola/00_tikz_legacy.tex}
\input{examples/legacy-nicola/13_highlighted_list_next_to_figure.tex}
\input{examples/legacy-nicola/14_animated_scenario_figure.tex}
).

\begin{frame}[plain,c]
    \centering
    {\Huge\bfseries \color{DarkBlue}01 Examples}\\[1em]
    {\large Ready-to-copy code for the techniques seen in Nicola's presentations}
\end{frame}

% ESEMPIO BASE: lista e figura TikZ sincronizzate dagli stessi overlay.
% Overlay 1: figura neutra; overlay 2, 3 e 4: evidenziazione di A, B e C.
% Overlay 5: A e C sono evidenziati insieme, mostrando un overlay non contiguo.
% La figura condivisa viene riusata anche dall'esempio advanced: il disegno
% resta lo stesso, ma gli intervalli passati ai suoi stili producono
% evidenziazioni diverse.

\begin{frame}
    \frametitle{02 Basic: highlighted list next to a figure}

    \newcommand{\stepNeutral}{1}
    \newcommand{\stepA}{2}
    \newcommand{\stepBOne}{3}
    \newcommand{\stepBTwo}{3}
    \newcommand{\stepC}{4}
    \newcommand{\stepAAndC}{5}

    \begin{columns}[T]
        \begin{column}{.4\textwidth}
            This list highlights the current step in the process:

            \vspace{.75em}

            \begin{itemize}\setlength{\itemsep}{1em}
                \item \alert<\stepA,\stepAAndC>{Step A}
                \item \alert<\stepBOne>{Step B1}
                \item \alert<\stepBTwo>{Step B2}
                \item \alert<\stepC,\stepAAndC>{Step C}
            \end{itemize}
        \end{column}
        \begin{column}{.6\textwidth}
            \begin{center}
                \input{examples/assets/figures/fig_overlay_graph.tex}
            \end{center}
        \end{column}
    \end{columns}
\end{frame}

% ESEMPIO BASE: definizione tecnica in un blocco Beamer.
% La struttura e la macro \angled sono riusate dall'esempio di Nicola;
% il contenuto resta neutro e sostituibile.

\begin{frame}
    \frametitle{03 Basic: defining a component}

    \begin{block}{Component definition}
        A \alert{component} is represented by
        \alert{$\angled{n, I, O}$}, where

        \vspace{.75em}

        \begin{itemize}\setlength{\itemsep}{.5em}
            \item \alert{$n$} is the component name;
            \item \alert{$I$} is the set of input elements;
            \item \alert{$O$} is the set of output elements.
        \end{itemize}
    \end{block}
\end{frame}

% ESEMPIO BASE: \pause per far comparire progressivamente i punti.
% L'immagine resta fissa mentre il testo a sinistra procede.

\begin{frame}
    \frametitle{04 Basic: progressive content with \texttt{\textbackslash pause}}

    \begin{columns}[T]
        \begin{column}{.65\textwidth}
            \begin{itemize}\setlength{\itemsep}{1em}
                \item First point, visible from the beginning.
                \pause
                \item Second point, revealed at the next overlay.
            \end{itemize}
        \end{column}
        \begin{column}{.35\textwidth}
            \includegraphics[width=.8\columnwidth]{assets/img/patrick_star_sad.png}
        \end{column}
    \end{columns}
\end{frame}

% ESEMPIO BASE: \only dentro overlayarea per sostituire il contenuto.
% L'overlayarea riserva lo spazio mentre cambia il blocco mostrato.

\begin{frame}
    \frametitle{05 Basic: replacing content with \texttt{\textbackslash only}}

    \begin{overlayarea}{\textwidth}{.65\textheight}
        \only<1>{
            \begin{block}{First view}
                This block is visible only on the first overlay.
            \end{block}
        }
        \only<2>{
            \begin{block}{Final view}
                The previous block is replaced by this content.
            \end{block}
        }
    \end{overlayarea}
\end{frame}

% ESEMPIO BASE: composizioni di tabelle riutilizzabili nelle slide.
% - due tabelle impilate con una didascalia per ciascuna;
% - intestazione su più righe con \multirow e \multicolumn;
% - tabella affiancata a un elenco puntato.
% booktabs migliora le righe della tabella; adjustbox evita che una tabella
% troppo larga esca dai margini della slide.

\begin{frame}
    \frametitle{06 Basic: stacked tables}

    \begin{adjustbox}{max totalsize={\textwidth}{.75\textheight},center}
        \begin{tabular}{c}
            \begin{tabular}{l c l}
                \toprule
                \textbf{\alert{Item}} & \textbf{\alert{Status}} & \textbf{\alert{Description}} \\
                \midrule
                Item A & ready & Short description \\
                Item B & pending & Another description \\
                \bottomrule
            \end{tabular}\\[-.1em]
            \textbf{First table}\\[1em]
            \begin{tabular}{p{.18\textwidth} c p{.3\textwidth}}
                \toprule
                \textbf{\alert{Item}} & \textbf{\alert{Value}} & \textbf{\alert{Comment}} \\
                \midrule
                Item C & 12 & A longer comment that wraps onto a second line \\
                Item D & 27 & Another comment with controlled cell width \\
                \bottomrule
            \end{tabular}\\[-.1em]
            \textbf{Second table}
        \end{tabular}
    \end{adjustbox}
\end{frame}

\begin{frame}
    \frametitle{06 Basic: table with grouped columns}

    \begin{columns}[T]
        \begin{column}{.62\textwidth}
            \begin{adjustbox}{max totalsize={\textwidth}{.7\textheight},center}
                \begin{tabular}{l r cc}
                    \toprule
                    \multirow{2}{*}{\textbf{\alert{Step}}}
                        & \multirow{2}{*}{\textbf{\alert{Value}}}
                        & \multicolumn{2}{c}{\alert{Cost}} \\
                    & & \textbf{\alert{Unit A}} & \textbf{\alert{Unit B}} \\
                    \midrule
                    Step 1 & 8 & 0.41 & 1.29 \\
                    Step 2 & 5 & 0.24 & 0.74 \\
                    Step 3 & 4 & 0.19 & 0.59 \\
                    \midrule
                    \textbf{Total} & \textbf{17} & \textbf{0.84} & \textbf{2.62} \\
                    \bottomrule
                \end{tabular}
            \end{adjustbox}
        \end{column}
        \begin{column}{.38\textwidth}
            \begin{itemize}\setlength{\itemsep}{.8em}
                \item The value is the measured quantity.
                \item Unit A and Unit B show two cost estimates.
                \item The final row gives the total.
            \end{itemize}
        \end{column}
    \end{columns}
\end{frame}

% ESEMPIO: grafici PGFPlots con dati esterni.
% La sorgente dei dati e' indipendente dal tipo di grafico:
% - grafici a linee letti da file .dat;
% - grafici a barre orizzontali letti da CSV.
% Le proporzioni, gli assi e le legende seguono gli esempi originali di Nicola.
% Il CSV e' invece contenuto dimostrativo neutro: le due configurazioni usano
% colonne diverse, non una copia degli stessi dati.

\begin{frame}
    \frametitle{07 Basic: completed sentence and line plots}

    \pgfmathsetmacro\numberTextStart{int(1)}
    \pgfmathsetmacro\numberTextEnd{int(\numberTextStart+0)}
    \pgfmathsetmacro\numberGraphStart{int(\numberTextEnd+1)}
    \pgfmathsetmacro\numberGraphEnd{int(\numberGraphStart+0)}

    \begin{overlayarea}{\textwidth}{.9\textheight}
        \only<\numberTextStart-\numberGraphEnd>{
            \alert{Measure} varying the number of
        }
        \only<\numberGraphStart-\numberGraphEnd>{
            \alert{parameter A} and \alert{parameter B}
        }

        \only<\numberTextStart-\numberTextEnd>{
            \vspace{1em}
            \begin{itemize}\setlength{\itemsep}{1em}
                \item \alert{parameter A} in $\{2,\ldots,5\}$
                \item \alert{parameter B} in $\{1,\ldots,5\}$
            \end{itemize}
        }

        \only<\numberGraphStart-\numberGraphEnd>{
            \vspace{1em}

            \pgfplotstableread{examples/assets/data/lines_01.dat}\dataTableParameterA
            \pgfplotstableread{examples/assets/data/lines_02.dat}\dataTableParameterB

            \pgfplotsset{every tick label/.append style={font=\normalsize}}

            \begin{adjustbox}{max totalsize={.99\textwidth}{.6\textheight},center}
                \begin{tabular}{cc}
                    \begin{tikzpicture}[scale=.5]
                        \begin{axis}[
                                xlabel={\normalsize \textbf{Parameter A}},
                                ylabel={\normalsize \textbf{Measure}},
                                xmin=2, xmax=5,
                                xtick={2, ..., 5},
                                ymin=4, ymax=11,
                                legend pos=north west,
                                legend columns=2,
                                legend style={/tikz/every even column/.append style={column sep=0.5cm}},
                                scale only axis
                            ]
                            \addplot[color=black, mark options={solid}]
                                table [x index=0, y index=1] {\dataTableParameterA};
                            \addlegendentry{1 unit}
                            \addplot[color=black, mark options={solid}, dashed]
                                table [x index=0, y index=2] {\dataTableParameterA};
                            \addlegendentry{2 units}
                            \addplot[color=black, mark options={solid}, dotted]
                                table [x index=0, y index=3] {\dataTableParameterA};
                            \addlegendentry{3 units}
                            \addplot[color=black, mark options={solid}, dashdotted]
                                table [x index=0, y index=4] {\dataTableParameterA};
                            \addlegendentry{4 units}
                            \addplot[color=black, mark options={solid},
                                dash pattern={on 1pt off 2pt on 1pt off 2pt on 3pt}]
                                table [x index=0, y index=5] {\dataTableParameterA};
                            \addlegendentry{5 units}
                        \end{axis}
                    \end{tikzpicture}
                    &
                    \begin{tikzpicture}[scale=.5]
                        \begin{axis}[
                                xlabel={\normalsize \textbf{Parameter B}},
                                xmin=1, xmax=5,
                                xtick={1,2,3,4,5},
                                ymin=4, ymax=11,
                                legend pos=north west,
                                scale only axis
                            ]
                            \addplot[color=black, mark options={solid}]
                                table [x index=0, y index=1] {\dataTableParameterB};
                            \addlegendentry{2 items}
                            \addplot[color=black, dashed, mark options={solid}]
                                table [x index=0, y index=2] {\dataTableParameterB};
                            \addlegendentry{3 items}
                            \addplot[color=black, dashdotted, mark options={solid}]
                                table [x index=0, y index=3] {\dataTableParameterB};
                            \addlegendentry{4 items}
                            \addplot[color=black, dotted, mark options={solid}]
                                table [x index=0, y index=4] {\dataTableParameterB};
                            \addlegendentry{5 items}
                        \end{axis}
                    \end{tikzpicture}
                    \\
                    {\footnotesize Varying parameter A} &
                    {\footnotesize Varying parameter B}
                \end{tabular}
            \end{adjustbox}

            \vspace{1em}
            $\Longrightarrow$ \alert{linear} increase of the measure
        }
    \end{overlayarea}
\end{frame}

\begin{frame}
    \frametitle{07 Advanced: horizontal bar plots from a CSV}

    \pgfplotstableread[col sep=comma]{examples/assets/data/data_comparison.csv}\dataComparison

    \begin{adjustbox}{max totalsize={.99\linewidth}{\textheight},center}
        \begin{tikzpicture}
            \begin{groupplot}[
                    group style={
                        group name=basic bar plots,
                        group size=2 by 1,
                        vertical sep=1.45cm,
                        xticklabels at=all,
                        yticklabels at=edge left,
                        xlabels at=edge bottom,
                        ylabels at=edge left
                    },
                    xlabel={Value},
                    ylabel={Category},
                    xticklabel style={font=\small},
                    yticklabel style={font=\small},
                    scaled y ticks=false,
                    scaled x ticks=false,
                    xbar,
                    ytick={1,2,3},
                    yticklabels={Group 1,Group 2,Group 3},
                    nodes near coords,
                    enlarge x limits={upper, value=0.18},
                    xmin=0,
                    width=.43\textwidth,
                    height=.48\textheight
                ]
                \nextgroupplot[title={\normalsize \textbf{Configuration 1}}]
                    \addplot[draw=black, fill=gray!50, bar width=4.5pt]
                        table[x=C1_Method_A, y expr=\coordindex+1]{\dataComparison};
                    \addplot[draw=black, fill=white, bar width=4.5pt]
                        table[x=C1_Method_B, y expr=\coordindex+1]{\dataComparison};
                    \coordinate (bar top left) at (rel axis cs:0,1);

                \nextgroupplot[
                        title={\normalsize \textbf{Configuration 2}},
                        legend style={
                            at={($(0,0)+(1cm,1cm)$)},
                            legend columns=2,
                            fill=none,
                            draw=none,
                            anchor=south,
                            inner sep=0.2em,
                            column sep=5pt,
                            align=center
                        },
                        legend to name=basicBarLegend
                    ]
                    \addplot[draw=black, fill=gray!50, bar width=4.5pt]
                        table[x=C2_Method_A, y expr=\coordindex+1]{\dataComparison};
                    \addlegendentry{Method A}
                    \addplot[draw=black, fill=white, bar width=4.5pt]
                        table[x=C2_Method_B, y expr=\coordindex+1]{\dataComparison};
                    \addlegendentry{Method B}
                    \coordinate (bar top right) at (rel axis cs:1,1);
            \end{groupplot}

            \coordinate (bar legend middle) at ($(bar top left)!.5!(bar top right)$);
            \coordinate (bar legend position)
                at (bar legend middle |- current bounding box.north);
            \node[above] at (bar legend position)
                {\pgfplotslegendfromname{basicBarLegend}};
        \end{tikzpicture}
    \end{adjustbox}
\end{frame}

% ESEMPIO AVANZATO: personaggi e fumetti riusati dalla tecnica di Nicola.
% Dimostra testo sostituito, due speaker simultanei e una scena finale,
% sincronizzati con l'highlighting del grafo.
% Il grafo incluso a sinistra e' la stessa figura condivisa dell'esempio
% basic: qui riceve intervalli diversi e quindi colorazioni diverse, sincronizzate
% con la comparsa dei personaggi e dei fumetti.

\begin{frame}
    \frametitle{08 Advanced: synchronized characters and speech}

    \newcommand{\stepA}{1}
    \newcommand{\stepBOne}{2}
    \newcommand{\stepBTwo}{3}
    \newcommand{\stepC}{4-5}
    \newcommand{\stepAAndC}{5}

    \begin{columns}[T]
        \begin{column}{.42\textwidth}
            \begin{center}
                \input{examples/assets/figures/fig_overlay_graph.tex}
            \end{center}
        \end{column}
        \begin{column}{.58\textwidth}
            \begin{center}
                \begin{tikzpicture}
                    \node[dave, minimum size=.8cm, visible on=<1-4>]
                        (dave-first) at (0,1.2) {};
                    \node[speaking question, callout absolute pointer=(dave-first.mouth),
                        anchor=west, visible on=<1>]
                        at ([xshift=1cm]dave-first.east) {\alert{Speaker 1:} Speech 1};
                    \node[speaking question, callout absolute pointer=(dave-first.mouth),
                        anchor=west, visible on=<2>]
                        at ([xshift=1cm]dave-first.east) {\alert{Speaker 1:} Speech 2};
                    \node[speaking question, callout absolute pointer=(dave-first.mouth),
                        anchor=west, visible on=<3>]
                        at ([xshift=1cm]dave-first.east) {\alert{Speaker 1:} Speech 3};
                    \node[speaking question, callout absolute pointer=(dave-first.mouth),
                        anchor=west, visible on=<4>]
                        at ([xshift=1cm]dave-first.east) {\alert{Speaker 1:} Speech 4};

                    \node[judge, minimum size=.8cm, visible on=<3-4>]
                        (judge-second) at (0,-1.2) {};
                    \node[speaking question, callout absolute pointer=(judge-second.mouth),
                        anchor=west, visible on=<3>]
                        at ([xshift=1cm]judge-second.east) {\alert{Speaker 2:} Speech 1};
                    \node[speaking question, callout absolute pointer=(judge-second.mouth),
                        anchor=west, visible on=<4>]
                        at ([xshift=1cm]judge-second.east) {\alert{Speaker 2:} Speech 2};

                    \node[dave, minimum size=.8cm, visible on=<5>]
                        (dave-final) at (0,0) {};
                    \node[speaking question, callout absolute pointer=(dave-final.mouth),
                        anchor=west, visible on=<5>]
                        at ([xshift=1cm]dave-final.east) {\alert{Final speaker:} Final speech};
                \end{tikzpicture}
            \end{center}
        \end{column}
    \end{columns}
\end{frame}

% ESEMPIO AVANZATO: riga di chevron con descrizioni e lista in primo piano.
% Riprende la struttura della tecnica originale di Nicola:
% - stile arrow per i chevron;
% - stile desc per le schede sotto ciascun chevron;
% - dissolvenza delle schede con text opacity e fill opacity;
% - lista in primo piano con visible on e highlighting progressivo.

\begin{frame}
    \frametitle{09 Advanced: chevrons and fading background}

    \pgfmathsetmacro\numberFadeStart{int(2)}
    \pgfmathsetmacro\numberItemA{int(\numberFadeStart+0)}
    \pgfmathsetmacro\numberItemB{int(\numberItemA+1)}
    \pgfmathsetmacro\numberItemC{int(\numberItemB+1)}
    \pgfmathsetmacro\numberFadeEnd{int(\numberItemC+0)}

    \begin{center}
        A sentence introducing the \alert{three phases}\ldots
    \end{center}

    \vspace{1em}

    \begin{adjustbox}{max totalsize={.9\textwidth}{\textheight},center}
        \input{examples/assets/figures/fig_chevrons.tex}
    \end{adjustbox}
\end{frame}

% ESEMPIO AVANZATO: componente parametrico TikZ con \pic.
% La definizione e le chiamate sono riusate direttamente da Nicola in
% examples/assets/figures/fig_pic.tex. Il componente riceve nome, posizione e overlay:
% la stessa costruzione puo' quindi essere istanziata piu' volte senza
% duplicare il disegno interno.

\begin{frame}
    \frametitle{10 Advanced: reusable component with \texttt{\textbackslash pic}}

    \begin{columns}[T]
        \begin{column}{.38\textwidth}
            The same component is instantiated twice:

            \vspace{.5em}

            \begin{itemize}\setlength{\itemsep}{.7em}
                \item \textbf{Name}: \texttt{a} or \texttt{b}; it becomes part of the node names.
                \item \textbf{Position}: the second argument moves the whole group.
                \item \textbf{Overlay}: the third argument selects when that group is highlighted.
                \item The relation uses the generated nodes
                    \texttt{box-a-3} and \texttt{box-b-3}.
            \end{itemize}
        \end{column}
        \begin{column}{.62\textwidth}
            \begin{adjustbox}{width=\linewidth, max height=.72\textheight, center}
                \input{examples/assets/figures/fig_pic.tex}
            \end{adjustbox}
        \end{column}
    \end{columns}
\end{frame}

% ESEMPIO AVANZATO: tre tecniche di callout nella stessa figura.
% La figura e i relativi stili sono riusati direttamente da Nicola in
% examples/assets/figures/fig_callouts.tex:
% 1. callout con codice, un puntatore e linguetta;
% 2. callout con più puntatori verso più nodi;
% 3. personaggio con fumetto.
% I tre casi compaiono su overlay successivi tramite visible on e on slide.

\begin{frame}
    \frametitle{11 Advanced: callouts and speech bubbles}

    \begin{adjustbox}{width=.95\linewidth, max height=.8\textheight, center}
        \input{examples/assets/figures/fig_callouts.tex}
    \end{adjustbox}
\end{frame}

% ESEMPIO AVANZATO: titolo con autori, loghi e finanziamenti.
% Il gruppo locale mantiene i metadati dell'esempio separati da quelli della
% presentazione principale.

\begingroup
\title{12 Advanced: title slide with logos and funding}
\author{\small \textbf{Speaker}\inst{1} \and Author Two\inst{1} \and Author Three\inst{2}}
\institute{\inst{1} Affiliation 1 \and \inst{2} Affiliation 2}
\date{}

\begin{frame}
    \titlepage

    \vfill

    \vspace{-2em}
    \begin{adjustbox}{max totalsize={\textwidth}{.125\textheight},center}
        \begin{tabular}{lr}
            \includegraphics[width=.05\textwidth]{assets/img/patrick_happy.jpg} &
            \includegraphics[width=.1\textwidth]{assets/img/patrick_scientist.jpg}
        \end{tabular}
    \end{adjustbox}

    \parMini{The work was partially supported by
        \emph{i)} project A, funded by the funder under the call number 1;
        \emph{ii)} project B, funded by another funder.
    }
\end{frame}
\endgroup


% Le due slide originali di Nicola hanno dipendenze specifiche, caricate
% solo quando queste slide legacy vengono incluse.
% Legacy Nicola: macro specifiche alla notazione e agli esperimenti originali.
% Gli helper generici sono definiti in 00_macro_common.tex e non vengono duplicati qui.

\newcommand*{\prop}{\ensuremath{p}}
\newcommand*{\target}{\ensuremath{ToC}}
\newcommand*{\targetSet}{\referred{\setStyle{TOC}}}
\newcommand*{\mechanism}{\ensuremath{m}}
% \newcommand*{\evidCollectionModel}{\ensuremath{\mathcal{E}}}
\newcommand*{\evidCollectionModel}{\ensuremath{E}}
\newcommand*{\evalFn}{\ensuremath{\mathcal{F}}}
% \newcommand*{\certModel}{\ensuremath{\mathcal{CM}}}
\newcommand*{\certModel}{\ensuremath{CM}}
\newcommand*{\certModelSet}{\referred{\setStyle{\certModel}}}
%\newcommand*{\cert}{\ensuremath{\mathcal{C}}}
\newcommand*{\cert}{\ensuremath{C}}
\newcommand*{\certSet}{\referred{\setStyle{\cert}}}
\newcommand*{\evid}{\ensuremath{ev}}
%\newcommand*{}{def}
\newcommand*{\attr}{\ensuremath{attr}}
\newcommand*{\repo}{\ensuremath{\mathcal{R}}}

\newcommand*{\signature}{\texttt{sig}}
\newcommand*{\signatureCA}[1][]{\signature\ensuremath{_{\text{CA}_{\text{#1}}}}}
\newcommand*{\signatureAL}[1][]{\signature\ensuremath{_{\text{AL}_{\text{#1}}}}}

\newcommand*{\dimension}{\ensuremath{d}}
\newcommand{\verificationStep}{\ensuremath{v}}

\newcommand{\targetInput}{\ensuremath{\mathcal{I}}}

%\newcommand{\dimension}{\ensuremath{d}}

\newcommand{\abbrvPreTraining}{\text{\textit{pt}}}
\newcommand{\abbrvAdaptation}{\text{\textit{ad}}}
\newcommand{\abbrvInference}{\text{\textit{in}}}
\newcommand{\abbrvArtifact}{\text{\textit{ar}}}

\newcommand{\viewPreTraining}{\ensuremath{\view_{\abbrvPreTraining}}}
\newcommand{\viewAdaptation}{\ensuremath{\view_{\abbrvAdaptation}}}
\newcommand{\viewInference}{\ensuremath{\view_{\abbrvInference}}}
\newcommand{\viewArtifact}{\ensuremath{\view_{\abbrvArtifact}}}
\newcommand{\viewAdaptationInference}{\ensuremath{\view_{\abbrvAdaptation+\abbrvInference}}}
\newcommand{\dimensionPreTraining}{\ensuremath{\dimension_{\abbrvPreTraining}}}
\newcommand{\dimensionAdaptation}{\ensuremath{\dimension_{\abbrvAdaptation}}}
\newcommand{\dimensionInference}{\ensuremath{\dimension_{\abbrvInference}}}
\newcommand{\dimensionAdaptationInference}{\ensuremath{\dimension_{\abbrvAdaptation+\abbrvInference}}}
\newcommand{\dimensionArtifact}{\ensuremath{\dimension_{\abbrvArtifact}}}

\newcommand{\succSymbol}{\ensuremath{\text{\textit{Succ}}}}
\newcommand{\succFn}[1]{\ensuremath{\succSymbol(#1)}}

\newcommand{\foundation}{\text{found}}

\newcommand{\testCaseOutput}{\ensuremath{out}}
\newcommand{\testCaseResult}{\ensuremath{res}}


%\newcommand{\hypergraphEvalRule}{\ensuremath{H_{\text{eval}}}}\
\newcommand{\hypergraphEvalRule}{\ensuremath{\mathcal{R}}}
\newcommand{\hypergraphAll}{\ensuremath{H_{\text{cert}}}}
\newcommand{\hypergraphInstance}{\ensuremath{\mathcal{H}}}
\newcommand{\verificationPath}{\ensuremath{\mathcal{V}}}

\newcommand{\hypergraphCertModel}{\ensuremath{H}}
%\newcommand{\hypergraphCert}{\ensuremath{H_{\text{cert}}}}
\NewDocumentCommand{\hypergraphCert}{o}{%
  \ensuremath{H_{\text{cert}\IfValueT{#1}{, #1}}}%
}

\newcommand{\hyperpath}{\ensuremath{\text{\emph{path}}}}

\newcommand{\FPathStarSymbol}{\ensuremath{\mathcal{F}}}
\newcommand{\BPathStarSymbol}{\ensuremath{\mathcal{B}}}
\newcommand{\BFPathStarSymbol}{\ensuremath{\mathcal{BF}}}
\newcommand{\FPathStarOf}[1]{\ensuremath{\FPathStarSymbol{(#1)}}}
\newcommand{\BPathStarOf}[1]{\ensuremath{\BPathStarSymbol{(#1)}}}
\newcommand{\BFPathStarOf}[1]{\ensuremath{\BFPathStarSymbol{(#1)}}}
\newcommand{\PathStarSymbol}{\ensuremath{\text{\emph{PathsSimple}}}}
\newcommand{\PathStarOf}[1]{\ensuremath{\PathStarSymbol{(#1)}}}

\newcommand{\evalNode}{\ensuremath{\evalFn}}
\newcommand{\evalNodeAll}{\ensuremath{\widehat{\evalFn}}}

% \newcommand{\nodes}{\ensuremath{V}}
% \newcommand{\edges}{\ensuremath{E}}
\newcommand{\nodes}{\ensuremath{N}}
\newcommand{\arcs}{\ensuremath{A}}

%\newcommand{\node}{\ensuremath{n}}

\newcommand{\modelName}[1]{\texttt{#1}}

\newcommand{\llmGemma}[1][short]{%
  \ifthenelse{\equal{#1}{short}}{%
    \modelName{gemma3:27b}%
  }{%
    \modelName{gemma3:27b-it-q4\_K\_M}%
  }%
}
\newcommand{\llmLlama}[1][short]{%
  \ifthenelse{\equal{#1}{short}}{%
    \modelName{llama3.3:70b}%
  }{%
    \modelName{llama3.3:70b-instruct-q5\_K\_M}%
  }%
}
\newcommand{\llmDeepseek}[1][short]{%
  \ifthenelse{\equal{#1}{short}}{%
    \modelName{deepseek-r1:70b}%
  }{%
    \modelName{deepseek-r1:70b}%
  }%
}
\newcommand{\llmMistral}[1][short]{%
  \ifthenelse{\equal{#1}{short}}{%
    \modelName{mistral:123b}%
  }{%
    \modelName{mistral-large:123b-instruct-2411-q5\_K\_M}%
  }%
}
\newcommand{\llmQwenCoder}[1][short]{%
  \ifthenelse{\equal{#1}{short}}{%
    \modelName{qwen3-coder:30b}%
  }{%
    \modelName{qwen3-coder:30b-a3b-q4\_K\_M}%
  }%
}
\newcommand{\llmGptOss}[1][short]{%
  \ifthenelse{\equal{#1}{short}}{%
    \modelName{gpt-oss:20b}%
  }{%
    \modelName{gpt-oss:20b}%
  }%
}
\newcommand{\llmCodeLlama}[1][short]{%
  \ifthenelse{\equal{#1}{short}}{%
    \modelName{codellama:70b}%
  }{%
    \modelName{codellama:70b-instruct-q5\_K\_M}%
  }%
}

% \newcommand{\llmGemma}{\modelName{gemma3:27b-it-q4\_K\_M}}
% \newcommand{\llmLlama}{\modelName{llama3.3:70b-instruct-q5\_K\_M}}
% \newcommand{\llmDeepseek}{\modelName{deepseek-r1:70b}}
% \newcommand{\llmMistral}{\modelName{mistral-large:123b-instruct-2411-q5\_K\_M}}
% \newcommand{\llmQwenCoder}{\modelName{qwen3-coder:30b-a3b-q4\_K\_M}}
% \newcommand{\llmGptOss}{\modelName{gpt-oss:20b}}
% \newcommand{\llmCodeLlama}{\modelName{codellama:70b-instruct-q5\_K\_M}}


\newcommand{\E}{\ensuremath{E}}
\newcommand{\Efirst}{\ensuremath{\E_{\text{first}}}}
\newcommand{\Elast}{\ensuremath{\E_{\text{last}}}}
\newcommand{\Erandom}{\ensuremath{\E_{\text{random}}}}
\newcommand{\Enull}{\ensuremath{\E_{\text{null}}}}

%\newcommand{\Prompts}{\ensuremath{\mathbb{P}}}

\newcommand{\llm}{\ensuremath{\text{\emph{app}}}}

\newcommand{\probe}{\ensuremath{pr}}
\newcommand{\myN}{\ensuremath{n}}
\newcommand{\myM}{\ensuremath{m}}
\newcommand{\accordance}{\ensuremath{\text{\emph{res}}}}
%\newcommand{\probeSet}{\ensuremath{P}}
\newcommand{\prompt}{\ensuremath{\text{\emph{prompt}}}}

\newcommand{\answer}{\ensuremath{\text{\emph{ans}}}}

\newcommand{\orderBias}{\ensuremath{\text{\emph{bias}}}}

\newlength{\oldtabcolsep}
\setlength{\oldtabcolsep}{\tabcolsep}

\newcommand{\tailSetSymbol}{\ensuremath{\text{\emph{tail}}}}
\newcommand{\headSetSymbol}{\ensuremath{\text{\emph{head}}}}

\newcommand{\node}{\ensuremath{\text{\emph{node}}}}
\newcommand{\arc}{\ensuremath{\text{\emph{arc}}}}

\newcommand{\headSet}[1]{\ensuremath{\headSetSymbol(#1)}}
\newcommand{\tailSet}[1]{\ensuremath{\tailSetSymbol(#1)}}

\newcommand{\guardFunctionSymbol}{\ensuremath{G}}
\newcommand{\guardFunction}[1]{\ensuremath{\guardFunctionSymbol(#1)}}
\newcommand{\evidenceFunctionSymbol}{\ensuremath{EV}}
\newcommand{\evidenceFunction}[1]{\ensuremath{\evidenceFunctionSymbol(#1)}}
\newcommand{\nil}{\texttt{nil}}
\newcommand{\error}{\texttt{error}}

%%%%%%%%%
% new macro

\newcommand{\propertyLabel}{\ensuremath{p}}
\newcommand{\behaviorSet}{\ensuremath{\mathcal{B}}}
\newcommand{\verificationStepTarget}{\ensuremath{t}}
\newcommand{\verificationStepActivity}{\ensuremath{s}}
\newcommand{\conditionGuardSymbol}{\ensuremath{\text{\emph{g}}}}
\newcommand{\conditionGuardFunction}[1]{\ensuremath{\conditionGuardSymbol(#1)}}

%\newcommand{\verificationPathAnnotated}{\ensuremath{\overline{\verificationPath}}}

\newcommand{\verificationGraphAnnotated}{\ensuremath{\overline{\behaviorSet}}}
\newcommand{\verificationNodesAnnotated}{\ensuremath{\overline{\nodes}}}
\newcommand{\verificationArcsAnnotated}{\ensuremath{\overline{\arcs}}}


\newcommand{\sourceFunctionSymbol}{\ensuremath{\text{\emph{source}}}}
\newcommand{\sourceFunction}[1]{\ensuremath{\sourceFunctionSymbol(#1)}}

\newcommand{\firstFunctionSymbol}{\ensuremath{\text{\emph{first}}}}
\newcommand{\firstFunction}[1]{\ensuremath{\firstFunctionSymbol(#1)}}

\newcommand{\lastFunctionSymbol}{\ensuremath{\text{\emph{last}}}}
\newcommand{\lastFunction}[1]{\ensuremath{\lastFunctionSymbol(#1)}}

\newcommand{\arcCounterSymbol}{\ensuremath{c}}
\newcommand{\arcCounterOf}[1]{\ensuremath{\arcCounterSymbol[#1]}}



\newcommand{\applTypeTextStandalone}{Type 1: Standalone}
\newcommand{\applTypeTextIntegratedArtifactGeneration}{Type 2: Integrated -- Artifact Generation}
\newcommand{\applTypeTextIntegratedDecisionMaking}{Type 3: Integrated -- Decision-Making}

\newcommand{\probeOutput}{\ensuremath{\text{\emph{out}}}}
\newcommand{\probeOutputTrue}{\ensuremath{\probeOutput^{\top}}}
\newcommand{\probeOutputFalse}{\ensuremath{\probeOutput^{\bot}}}

\newcommand{\FKGLSymbolFunction}{\ensuremath{\text{\emph{FK}}}}
\newcommand{\FKGLFunction}[1]{\ensuremath{\FKGLSymbolFunction(#1)}}
\newcommand{\DLSymbolFunction}{\ensuremath{\text{\emph{DL}}}}
\newcommand{\DLFunction}[1]{\ensuremath{\DLSymbolFunction(#1)}}

\newcommand{\arcsCardinality}{\ensuremath{\vert \arcs \vert}}
\newcommand{\nodesCardinality}{\ensuremath{\vert \nodes \vert}}

\newcommand{\nodesFailed}{\ensuremath{\overline{F}}}


\newcommand{\fullCiteMini}[1]{\parMini{\fullcite{#1}}}


\newcommand{\fullCiteAlmostMini}[1]{\parAlmostMini{\fullcite{#1}}}

\newcommand*{\service}[1]{\ensuremath{s_{#1}}}

\newcommand*{\pavail}{\ensuremath{\prop_\text{avail}}}
\newcommand*{\ppriv}{\ensuremath{\prop_{\text{priv}}}}
\newcommand*{\pConfTransit}{\ensuremath{\prop_{\text{conf\_net}}}}
\newcommand*{\pConfStorage}{\ensuremath{\prop_{\text{conf\_str}}}}


%\newcommand*{\relwordTargetCloudTrad}{def}
\newcommand*{\abbrvServiceStorage}{\text{str}}
\newcommand*{\abbrvServiceDatabase}{\text{db}}
\newcommand*{\abbrvServiceAuth}{\text{auth}}
\newcommand*{\abbrvServiceNetwork}{\text{net}}
\newcommand*{\abbrvServiceKube}{\text{k8s}}
\newcommand*{\abbrvServiceAnalytics}{\text{dawp}}
% \newcommand*{\abbrvServiceNetworkInt}{\text{netinf}}
% \newcommand*{\abbrvServiceStorageInt}{\text{strinf}}

\newcommand*{\abbrvServiceNetworkIntCloud}{\text{netinf\_cl}}
\newcommand*{\abbrvServiceStorageIntCloud}{\text{strinf\_cl}}
\newcommand*{\abbrvServiceNetworkIntEdge}{\text{netinf\_edge}}
\newcommand*{\abbrvServiceStorageIntEdge}{\text{strinf\_edge}}

\newcommand*{\abbrvMechanismUserAuth}{\text{user}}
\newcommand*{\abbrvMechanismScheduling}{\text{sched}}
\newcommand*{\abbrvMechanismData}{\text{data}}

\newcommand*{\abbrvHandlerApp}{\text{app}}
\newcommand*{\abbrvHandlerCSP}{\text{csp}}
\newcommand*{\abbrvHandlerISP}{\text{isp}}

\newcommand*{\CA}{\emph{\gls{CA}}}
\newcommand*{\CAApplication}{\CA\ensuremath{_{\abbrvHandlerApp}}}
\newcommand*{\CACSP}{\CA\ensuremath{_{\abbrvHandlerCSP}}}
\newcommand*{\CAISP}{\CA\ensuremath{_{\abbrvHandlerISP}}}

\newcommand*{\AL}{\emph{\gls{AL}}}
\newcommand*{\ALApplication}{\AL\ensuremath{_{\abbrvHandlerApp}}}
\newcommand*{\ALCSP}{\AL\ensuremath{_{\abbrvHandlerCSP}}}
\newcommand*{\ALISP}{\AL\ensuremath{_{\abbrvHandlerISP}}}

\newcommand{\approachSOTA}{\text{\emph{SOTA}}}
\newcommand{\approachOur}{\text{\emph{multi}}}
\newcommand{\certTimeSymbol}{\ensuremath{T}}
\newcommand{\certTimeFunction}[1]{\ensuremath{\certTimeSymbol}(#1)}

\newcommand*{\certStatusStyle}[1]{\texttt{#1}}
\newcommand*{\certStatusValid}[1]{\certStatusStyle{valid(}#1\certStatusStyle{)}} % chktex 9 chktex 36
\newcommand*{\certStatusSuspended}[1]{\certStatusStyle{suspended(}#1\certStatusStyle{)}} % chktex 9 chktex 36
\newcommand*{\certStatusRevoked}[1]{\certStatusStyle{revoked(}#1\certStatusStyle{)}} % chktex 9 chktex 36

% Legacy Nicola: colori, helper e stili specifici delle figure originali.
% Gli stili generici sono definiti in 00_tikz_common.tex e non vengono duplicati qui.

\colorlet{PoisonedRedForText}{white!70!red!100}
\colorlet{PoisonedRedOut}{red!70}
\colorlet{UnPoisonedGreen}{green!40!black!40}
\colorlet{UnPoisonedGreenForText}{green!30!white!90}
\definecolor{NotSoDarkGrey}{RGB}{160, 160, 160}
\definecolor{Completed}{RGB}{82,163,0}
\definecolor{InProgress}{RGB}{204, 102, 0}
\definecolor{colorBlindBiOrange}{HTML}{E66100}
\definecolor{colorBlindBiViolet}{HTML}{5D3A9B}

\newcommand{\normalwidth}{0.3pt}
\newcommand{\thickwidth}{0.8pt}
\newcommand{\forhash}{for hash}
\newcommand{\forrisk}{for risk}
\newcommand{\forToolTraining}{for tool training}
\newcommand{\forToolTesting}{for tool testing}
\newcommand{\poisonedRF}{1}
\newcommand{\headerEnsemblePoisoned}{Hash\-/Based Ensemble (Poisoned Training Sets)}
\newcommand{\headerEnsembleNonPoisoned}{Hash\-/Based Ensemble (Clean Training Sets)}
\newcommand{\headerEnsembleDefault}{Hash\-/Based Ensemble}
\newcommand{\headerRfTestingPoisoned}{Poisoned}
\newcommand{\headerRfTestingNonPoisoned}{Non\-/Poisoned}
\newcommand{\headerRfOther}{}
\newcommand{\codePoisoned}{Poisoned}
\newcommand{\codeNonPoisoned}{Non\-/Poisoned}
\newcommand{\devilLarge}{Large}
\newcommand{\devilSmall}{Small}

\newcommand{\includegraphicswithwisibility}[3]{
  \begin{tikzpicture}
    \node[visible with opacity on={#1}] at (0,0) {\includegraphics[#2]{#3}};
  \end{tikzpicture}
}

\tikzset{
  dataset/.style={
    draw,
    cylinder,
    shape border rotate=90,
    align=center,
  },
  rf/.style={
    draw,
    rounded corners,
    fill=white,
  },
  test path/.style={dashed},
  connections/.style={very thin},
  partition/.style={
    draw,
    cylinder,
    shape border rotate=90,
    align=center,
    minimum height=.6cm,
    minimum width=.35cm,
    fill=white,
  },
  poisoning link/.style={
    link,
    decorate,
    decoration={
      amplitude=.5pt,
      segment length=4pt,
      zigzag,
    },
    ultra thin,
  },
  data point/.style={
    draw,
    rounded corners,
    font={\tiny},
    fill=white,
  },
  data point moving/.style={
    circle,
    draw,
    fill=black,
    inner sep=1pt,
  },
  poisoned/.style={fill=PoisonedRedOut},
  majority voting/.style={
    isosceles triangle,
    draw,
  },
  highlight path/.style={
    highlight base,
    draw=LightGoldenRod,
  },
  highlight node/.style={
    highlight base,
    fill=HighlightFill,
  },
  highlight double/.style={
    double=LightGoldenRod,
    double distance=1.2pt,
    line width=.3pt,
  },
  arrow path unpoisoned/.append style={
    double=UnPoisonedGreen,
    arrows={-Latex[fill=UnPoisonedGreen]},
  },
  arrow path poisoned/.append style={
    double=PoisonedRedOut,
    arrows={-Latex[fill=PoisonedRedOut]},
  },
  highlight path poisoned/.style={},
  highlight path unpoisoned/.style={},
  invisible with opacity/.style={opacity=0.025},
  visible with opacity on/.style={alt=#1{}{invisible with opacity}},
  highlight line in pgfplot/.style={
    LightGoldenRod,
    line width=3pt,
  },
  poisoning attack/.style={
    draw,
    fill=white,
  },
  super container/.style={
    draw,
    dashed,
  },
  low link/.style={
    link,
    ultra thin,
  },
  bar/.style={
    thin,
    dotted,
  },
  from test/.style={ultra thin},
  highlight accuracy/.style={ultra thick},
  explanation/.style={font={\Large}},
  explanation link/.style={>=Latex},
  line am/.style={},
  line m2/.style={dashed},
  service in certification process/.style={
    chamfered rectangle,
    chamfered rectangle angle=30,
    draw,
  },
  connector assurance/.style={
    >=Stealth,
    dashed,
  },
  probe link/.style={
    link,
    dashdotted,
  },
  ca cert/.style={
    entity,
    inner sep=0pt,
  },
  step separator/.style={
    draw,
    thin,
    dotted,
  },
  container line/.style={
    very thin,
    rounded corners,
  },
  label desc/.style={},
  label desc with shape/.style={
    label desc,
    rounded corners,
    draw,
  },
  lattice layer label/.style={
    text width=2.9cm,
    font={\footnotesize},
  },
  not chosen color/.style={black!50},
  lattice layer separator/.style={dashed},
  service label/.style={
    fill=white,
    circle,
    draw,
    very thick,
    inner sep=2.5pt,
  },
  lattice element bound/.style={
    fill=black!15,
    double,
  },
  lattice element selected/.style={
    fill=Goldenrod!15,
    double,
  },
  lattice element not chosen/.style={
    not chosen color,
    densely dotted,
  },
  lattice element not chosen link/.style={
    not chosen color,
    densely dotted,
  },
  service element not chosen/.style={
    text=black!50,
    draw=black!50,
    dashed,
  },
  aggregation block/.style={
    circle,
    draw,
    inner sep=1pt,
  },
  link comment/.style={dashed},
  basic graph node/.style={
    circle,
    draw,
    minimum width=.6cm,
  },
  verification step/.style={
    basic graph node,
    fill=white,
  },
  evaluation function/.style={
    basic graph node,
    double,
    fill=white,
  },
  evaluation function global/.style={
    draw,
    double,
    rectangle,
    minimum width=.6cm,
    minimum height=.6cm,
  },
  hyperarc/.style={>=Latex},
  pics/hyperarc between fan in/.style n args={5}{
    code={
      \path let \p1 = ($(#1)!.5!(#2)$), \p2 = ($(#2)!.5!(#3)$) in
        coordinate (#4) at (\x1, \y2);
      \foreach \src in {#1, #2} {
        \draw[-, hyperarc, #5] (\src) -- (#4);
      }
      \draw[->, hyperarc, #5] (#4) -- (#3);
    }
  },
  pics/hyperarc between fan out/.style n args={5}{
    code={
      \path let \p1 = ($(#1)!.5!(#2)$), \p2 = ($(#2)!.5!(#3)$) in
        coordinate (#4) at (\x1, \y2);
      \foreach \src in {#1, #2} {
        \draw[<-, hyperarc, #5] (\src) -- (#4);
      }
      \draw[-, hyperarc, #5] (#4) -- (#3);
    }
  },
  visited node/.style={},
  visited arc/.style={thick},
  pruned node/.style={fill=gray!30},
  pruned arc/.style={draw=gray!50},
  evidence content/.style={font=\scriptsize},
  orange in violet out double/.style={
    fill=colorBlindBiOrange,
    double=colorBlindBiViolet,
    draw=colorBlindBiViolet,
    very thick,
  },
  orange out/.style={
    draw=colorBlindBiOrange,
    very thick,
  },
  highlighted hyperarc in path/.style={thick},
  hyperarc in verification path/.style={
    very thick,
    highlight path,
  },
  node in verification path/.style={
    very thick,
    highlight path,
  },
  certificate callout/.style={
    rectangle callout,
    fill=VeryLightGrey!50,
    draw,
  },
}

% !TEX root = ../main.tex
% ESEMPIO 0a: slide Background con elenco evidenziato passo passo accanto a una figura.
%   Ogni punto dell'elenco usa \alert<n>, che lo evidenzia solo alla sovrapposizione (overlay) n. Gli stessi
%   numeri sono usati dalla figura (assets/tex/01_fig_cert_process.tex) per evidenziare la parte corrispondente.
%   Da copiare: le macro dei numeri, le due colonne e l'\input della figura dentro adjustbox.

\begin{frame}
    \frametitle{13 Legacy Nicola: highlighted list next to a figure}
    % GUIDA: il concetto, lo schema o il contesto di cui parla il lavoro, con i suoi elementi. Meglio una figura che una definizione
    % lunga: ogni elemento si evidenzia (\alert<n>) mentre lo nomini.

    % Numero di sovrapposizione (overlay) in cui compare ogni passo. Il primo è 2 perché alla
    % sovrapposizione 1 si vede la figura senza evidenziazioni; ogni passo parte dal precedente.
    % Gli stessi nomi sono usati da assets/tex/01_fig_cert_process.tex per evidenziare la figura.
    \pgfmathsetmacro\numberProperty{int(2)}
    \pgfmathsetmacro\numberTarget{int(\numberProperty+1)}
    \pgfmathsetmacro\numberEvidenceCollectionModel{int(\numberTarget+1)}
    \pgfmathsetmacro\numberCertificationModel{int(\numberEvidenceCollectionModel+1)}
    \pgfmathsetmacro\numberExecution{int(\numberCertificationModel+1)}
    \pgfmathsetmacro\numberCertificate{int(\numberExecution+1)}

    \begin{columns}[T]
        \begin{column}{.45\textwidth}

            CHANGE: one sentence introducing the \alert{concept} and its \alert{key words}, followed by the elements

            \vspace{.5em}

            % Un punto per ogni passo: \alert<n> lo evidenzia solo alla sovrapposizione n.
            \begin{itemize}\setlength\itemsep{1em}
                \item \alert<\numberProperty>{CHANGE: step 1}
                \item \alert<\numberTarget>{CHANGE: step 2}
                \item \alert<\numberEvidenceCollectionModel>{CHANGE: step 3}
                \item \alert<\numberCertificationModel>{CHANGE: step 4}
                \item \alert<\numberExecution>{CHANGE: step 5}
                \item \alert<\numberCertificate>{CHANGE: step 6}
            \end{itemize}

        \end{column}
        \begin{column}{.55\textwidth}
            \begin{adjustbox}{max totalsize={\linewidth}{.8\textheight},center}
                \input{assets/tex/01_fig_cert_process.tex}
            \end{adjustbox}
        \end{column}
    \end{columns}

\end{frame}

% !TEX root = ../main.tex
% ESEMPIO 0b: slide Motivations con una figura animata (personaggi e bolle di dialogo).
%   La figura può essere un'immagine qualsiasi. Qui i numeri di sovrapposizione (overlay) sono definiti nel frame:
%   ogni passo parte dal precedente (+1) o resta uguale (+0) quando due cose compaiono insieme.
%   Gli stessi nomi sono usati da assets/tex/02_fig_cert_process.tex.
%   Da copiare: la struttura (macro dei numeri nel frame, figura in assets/tex/, \input dentro adjustbox).
%   La stessa tecnica si riusa nell'approccio, per mostrare una figura con i colori che evidenziano i passi.

\begin{frame}
    \frametitle{14 Legacy Nicola: animated scenario figure}
    % GUIDA: il problema, e perché lo stato dell'arte non basta. Meglio mostrarlo su un esempio concreto (una figura di uno scenario),
    % che può tornare più avanti nell'approccio.

    \pgfmathsetmacro\numberStart{int(1)}
    % passo: si evidenzia l'elemento principale e le sue dipendenze
    \pgfmathsetmacro\numberServiceTopStart{int(\numberStart+0)}
    \pgfmathsetmacro\numberHighlightDependencyStart{int(\numberServiceTopStart+0)}
    % passo: nuova bolla sull'elemento principale, che poi smette di essere evidenziato
    \pgfmathsetmacro\numberServiceTopEvidStart{int(\numberHighlightDependencyStart+1)}
    \pgfmathsetmacro\numberServiceTopEvidEnd{int(\numberServiceTopEvidStart+0)}
    \pgfmathsetmacro\numberServiceTopEnd{int(\numberServiceTopEvidEnd+0)}
    % passo: restano evidenziati gli altri elementi, nuova bolla
    \pgfmathsetmacro\numberCalloutFocusOnOtherStart{int(\numberServiceTopEvidEnd+1)}
    \pgfmathsetmacro\numberCalloutFocusOnOtherEnd{int(\numberCalloutFocusOnOtherStart+0)}
    % passo: doppio evidenziamento di più elementi
    \pgfmathsetmacro\numberHighlightDependencyEnd{int(\numberCalloutFocusOnOtherEnd+0)}
    \pgfmathsetmacro\numberHighlightDependencySomeOkStart{int(\numberCalloutFocusOnOtherStart+0)}
    \pgfmathsetmacro\numberHighlightDependencySomeOkEnd{int(\numberHighlightDependencySomeOkStart+0)}
    % passo: doppio evidenziamento di un altro elemento, nuova bolla
    \pgfmathsetmacro\numberHighlightDependencySomeNotOkStart{int(\numberHighlightDependencySomeOkEnd+1)}
    \pgfmathsetmacro\numberHighlightDependencySomeNotOkStartPartOne{int(\numberHighlightDependencySomeNotOkStart+0)}
    \pgfmathsetmacro\numberHighlightDependencySomeNotOkStartPartOneEnd{int(\numberHighlightDependencySomeNotOkStartPartOne+0)}
    \pgfmathsetmacro\numberHighlightDependencySomeNotOkStartPartTwoStart{int(\numberHighlightDependencySomeNotOkStartPartOneEnd+0)}
    \pgfmathsetmacro\numberHighlightDependencySomeNotOkStartPartTwoEnd{int(\numberHighlightDependencySomeNotOkStartPartTwoStart+0)}
    \pgfmathsetmacro\numberHighlightDependencySomeNotOkEnd{int(\numberHighlightDependencySomeNotOkStartPartTwoEnd+0)}
    \pgfmathsetmacro\numberEnd{int(\numberHighlightDependencySomeNotOkEnd+0)}
    % passo finale: ogni elemento mostra la sua bolla e compare l'immagine
    \pgfmathsetmacro\numberCertStart{int(\numberHighlightDependencySomeNotOkEnd+1)}

    \pgfmathsetmacro\numberPatrickHappy{int(\numberCertStart+1)}
    \pgfmathsetmacro\numberCertEnd{int(\numberPatrickHappy+0)}

    \begin{adjustbox}{max totalsize={.9\linewidth}{\textheight},center}
        \input{assets/tex/02_fig_cert_process.tex}
    \end{adjustbox}

\end{frame}

).

\begin{frame}[plain,c]
    \centering
    {\Huge\bfseries \color{DarkBlue}01 Examples}\\[1em]
    {\large Ready-to-copy code for the techniques seen in Nicola's presentations}
\end{frame}

% ESEMPIO BASE: lista e figura TikZ sincronizzate dagli stessi overlay.
% Overlay 1: figura neutra; overlay 2, 3 e 4: evidenziazione di A, B e C.
% Overlay 5: A e C sono evidenziati insieme, mostrando un overlay non contiguo.
% La figura condivisa viene riusata anche dall'esempio advanced: il disegno
% resta lo stesso, ma gli intervalli passati ai suoi stili producono
% evidenziazioni diverse.

\begin{frame}
    \frametitle{02 Basic: highlighted list next to a figure}

    \newcommand{\stepNeutral}{1}
    \newcommand{\stepA}{2}
    \newcommand{\stepBOne}{3}
    \newcommand{\stepBTwo}{3}
    \newcommand{\stepC}{4}
    \newcommand{\stepAAndC}{5}

    \begin{columns}[T]
        \begin{column}{.4\textwidth}
            This list highlights the current step in the process:

            \vspace{.75em}

            \begin{itemize}\setlength{\itemsep}{1em}
                \item \alert<\stepA,\stepAAndC>{Step A}
                \item \alert<\stepBOne>{Step B1}
                \item \alert<\stepBTwo>{Step B2}
                \item \alert<\stepC,\stepAAndC>{Step C}
            \end{itemize}
        \end{column}
        \begin{column}{.6\textwidth}
            \begin{center}
                \begin{tikzpicture}[node distance=2.5em and 2.8em]
    \node[draw, rounded corners, fill=white,
        on slide={<\stepA,\stepAAndC> highlight light}] (a) {Step A};
    \node[draw, rounded corners, fill=white, below left=of a,
        on slide={<\stepBOne> highlight light}] (b1) {Step B1};
    \node[draw, rounded corners, fill=white, below right=of a,
        on slide={<\stepBTwo> highlight light}] (b2) {Step B2};
    \node[draw, rounded corners, fill=white, below=of b1,
        on slide={<\stepC,\stepAAndC> highlight light}] (c) {Step C};

    \draw[->, link,
        on slide={<\stepA,\stepAAndC> highlight link}]
        (a) -- node[on link, on slide={<\stepA,\stepAAndC> highlight text}] {branch 1} (b1);
    \draw[->, link,
        on slide={<\stepA,\stepAAndC> highlight link}]
        (a) -- node[on link, on slide={<\stepA,\stepAAndC> highlight text}] {branch 2} (b2);
    \draw[->, link,
        on slide={<\stepBOne> highlight link}]
        (b1) -- node[on link, on slide={<\stepBOne> highlight text}] {merge} (c);
    \draw[->, link,
        on slide={<\stepBTwo> highlight link}]
        (b2) -- node[on link, on slide={<\stepBTwo> highlight text}] {merge} (c);
\end{tikzpicture}

            \end{center}
        \end{column}
    \end{columns}
\end{frame}

% ESEMPIO BASE: definizione tecnica in un blocco Beamer.
% La struttura e la macro \angled sono riusate dall'esempio di Nicola;
% il contenuto resta neutro e sostituibile.

\begin{frame}
    \frametitle{03 Basic: defining a component}

    \begin{block}{Component definition}
        A \alert{component} is represented by
        \alert{$\angled{n, I, O}$}, where

        \vspace{.75em}

        \begin{itemize}\setlength{\itemsep}{.5em}
            \item \alert{$n$} is the component name;
            \item \alert{$I$} is the set of input elements;
            \item \alert{$O$} is the set of output elements.
        \end{itemize}
    \end{block}
\end{frame}

% ESEMPIO BASE: \pause per far comparire progressivamente i punti.
% L'immagine resta fissa mentre il testo a sinistra procede.

\begin{frame}
    \frametitle{04 Basic: progressive content with \texttt{\textbackslash pause}}

    \begin{columns}[T]
        \begin{column}{.65\textwidth}
            \begin{itemize}\setlength{\itemsep}{1em}
                \item First point, visible from the beginning.
                \pause
                \item Second point, revealed at the next overlay.
            \end{itemize}
        \end{column}
        \begin{column}{.35\textwidth}
            \includegraphics[width=.8\columnwidth]{assets/img/patrick_star_sad.png}
        \end{column}
    \end{columns}
\end{frame}

% ESEMPIO BASE: \only dentro overlayarea per sostituire il contenuto.
% L'overlayarea riserva lo spazio mentre cambia il blocco mostrato.

\begin{frame}
    \frametitle{05 Basic: replacing content with \texttt{\textbackslash only}}

    \begin{overlayarea}{\textwidth}{.65\textheight}
        \only<1>{
            \begin{block}{First view}
                This block is visible only on the first overlay.
            \end{block}
        }
        \only<2>{
            \begin{block}{Final view}
                The previous block is replaced by this content.
            \end{block}
        }
    \end{overlayarea}
\end{frame}

% ESEMPIO BASE: composizioni di tabelle riutilizzabili nelle slide.
% - due tabelle impilate con una didascalia per ciascuna;
% - intestazione su più righe con \multirow e \multicolumn;
% - tabella affiancata a un elenco puntato.
% booktabs migliora le righe della tabella; adjustbox evita che una tabella
% troppo larga esca dai margini della slide.

\begin{frame}
    \frametitle{06 Basic: stacked tables}

    \begin{adjustbox}{max totalsize={\textwidth}{.75\textheight},center}
        \begin{tabular}{c}
            \begin{tabular}{l c l}
                \toprule
                \textbf{\alert{Item}} & \textbf{\alert{Status}} & \textbf{\alert{Description}} \\
                \midrule
                Item A & ready & Short description \\
                Item B & pending & Another description \\
                \bottomrule
            \end{tabular}\\[-.1em]
            \textbf{First table}\\[1em]
            \begin{tabular}{p{.18\textwidth} c p{.3\textwidth}}
                \toprule
                \textbf{\alert{Item}} & \textbf{\alert{Value}} & \textbf{\alert{Comment}} \\
                \midrule
                Item C & 12 & A longer comment that wraps onto a second line \\
                Item D & 27 & Another comment with controlled cell width \\
                \bottomrule
            \end{tabular}\\[-.1em]
            \textbf{Second table}
        \end{tabular}
    \end{adjustbox}
\end{frame}

\begin{frame}
    \frametitle{06 Basic: table with grouped columns}

    \begin{columns}[T]
        \begin{column}{.62\textwidth}
            \begin{adjustbox}{max totalsize={\textwidth}{.7\textheight},center}
                \begin{tabular}{l r cc}
                    \toprule
                    \multirow{2}{*}{\textbf{\alert{Step}}}
                        & \multirow{2}{*}{\textbf{\alert{Value}}}
                        & \multicolumn{2}{c}{\alert{Cost}} \\
                    & & \textbf{\alert{Unit A}} & \textbf{\alert{Unit B}} \\
                    \midrule
                    Step 1 & 8 & 0.41 & 1.29 \\
                    Step 2 & 5 & 0.24 & 0.74 \\
                    Step 3 & 4 & 0.19 & 0.59 \\
                    \midrule
                    \textbf{Total} & \textbf{17} & \textbf{0.84} & \textbf{2.62} \\
                    \bottomrule
                \end{tabular}
            \end{adjustbox}
        \end{column}
        \begin{column}{.38\textwidth}
            \begin{itemize}\setlength{\itemsep}{.8em}
                \item The value is the measured quantity.
                \item Unit A and Unit B show two cost estimates.
                \item The final row gives the total.
            \end{itemize}
        \end{column}
    \end{columns}
\end{frame}

% ESEMPIO: grafici PGFPlots con dati esterni.
% La sorgente dei dati e' indipendente dal tipo di grafico:
% - grafici a linee letti da file .dat;
% - grafici a barre orizzontali letti da CSV.
% Le proporzioni, gli assi e le legende seguono gli esempi originali di Nicola.
% Il CSV e' invece contenuto dimostrativo neutro: le due configurazioni usano
% colonne diverse, non una copia degli stessi dati.

\begin{frame}
    \frametitle{07 Basic: completed sentence and line plots}

    \pgfmathsetmacro\numberTextStart{int(1)}
    \pgfmathsetmacro\numberTextEnd{int(\numberTextStart+0)}
    \pgfmathsetmacro\numberGraphStart{int(\numberTextEnd+1)}
    \pgfmathsetmacro\numberGraphEnd{int(\numberGraphStart+0)}

    \begin{overlayarea}{\textwidth}{.9\textheight}
        \only<\numberTextStart-\numberGraphEnd>{
            \alert{Measure} varying the number of
        }
        \only<\numberGraphStart-\numberGraphEnd>{
            \alert{parameter A} and \alert{parameter B}
        }

        \only<\numberTextStart-\numberTextEnd>{
            \vspace{1em}
            \begin{itemize}\setlength{\itemsep}{1em}
                \item \alert{parameter A} in $\{2,\ldots,5\}$
                \item \alert{parameter B} in $\{1,\ldots,5\}$
            \end{itemize}
        }

        \only<\numberGraphStart-\numberGraphEnd>{
            \vspace{1em}

            \pgfplotstableread{examples/assets/data/lines_01.dat}\dataTableParameterA
            \pgfplotstableread{examples/assets/data/lines_02.dat}\dataTableParameterB

            \pgfplotsset{every tick label/.append style={font=\normalsize}}

            \begin{adjustbox}{max totalsize={.99\textwidth}{.6\textheight},center}
                \begin{tabular}{cc}
                    \begin{tikzpicture}[scale=.5]
                        \begin{axis}[
                                xlabel={\normalsize \textbf{Parameter A}},
                                ylabel={\normalsize \textbf{Measure}},
                                xmin=2, xmax=5,
                                xtick={2, ..., 5},
                                ymin=4, ymax=11,
                                legend pos=north west,
                                legend columns=2,
                                legend style={/tikz/every even column/.append style={column sep=0.5cm}},
                                scale only axis
                            ]
                            \addplot[color=black, mark options={solid}]
                                table [x index=0, y index=1] {\dataTableParameterA};
                            \addlegendentry{1 unit}
                            \addplot[color=black, mark options={solid}, dashed]
                                table [x index=0, y index=2] {\dataTableParameterA};
                            \addlegendentry{2 units}
                            \addplot[color=black, mark options={solid}, dotted]
                                table [x index=0, y index=3] {\dataTableParameterA};
                            \addlegendentry{3 units}
                            \addplot[color=black, mark options={solid}, dashdotted]
                                table [x index=0, y index=4] {\dataTableParameterA};
                            \addlegendentry{4 units}
                            \addplot[color=black, mark options={solid},
                                dash pattern={on 1pt off 2pt on 1pt off 2pt on 3pt}]
                                table [x index=0, y index=5] {\dataTableParameterA};
                            \addlegendentry{5 units}
                        \end{axis}
                    \end{tikzpicture}
                    &
                    \begin{tikzpicture}[scale=.5]
                        \begin{axis}[
                                xlabel={\normalsize \textbf{Parameter B}},
                                xmin=1, xmax=5,
                                xtick={1,2,3,4,5},
                                ymin=4, ymax=11,
                                legend pos=north west,
                                scale only axis
                            ]
                            \addplot[color=black, mark options={solid}]
                                table [x index=0, y index=1] {\dataTableParameterB};
                            \addlegendentry{2 items}
                            \addplot[color=black, dashed, mark options={solid}]
                                table [x index=0, y index=2] {\dataTableParameterB};
                            \addlegendentry{3 items}
                            \addplot[color=black, dashdotted, mark options={solid}]
                                table [x index=0, y index=3] {\dataTableParameterB};
                            \addlegendentry{4 items}
                            \addplot[color=black, dotted, mark options={solid}]
                                table [x index=0, y index=4] {\dataTableParameterB};
                            \addlegendentry{5 items}
                        \end{axis}
                    \end{tikzpicture}
                    \\
                    {\footnotesize Varying parameter A} &
                    {\footnotesize Varying parameter B}
                \end{tabular}
            \end{adjustbox}

            \vspace{1em}
            $\Longrightarrow$ \alert{linear} increase of the measure
        }
    \end{overlayarea}
\end{frame}

\begin{frame}
    \frametitle{07 Advanced: horizontal bar plots from a CSV}

    \pgfplotstableread[col sep=comma]{examples/assets/data/data_comparison.csv}\dataComparison

    \begin{adjustbox}{max totalsize={.99\linewidth}{\textheight},center}
        \begin{tikzpicture}
            \begin{groupplot}[
                    group style={
                        group name=basic bar plots,
                        group size=2 by 1,
                        vertical sep=1.45cm,
                        xticklabels at=all,
                        yticklabels at=edge left,
                        xlabels at=edge bottom,
                        ylabels at=edge left
                    },
                    xlabel={Value},
                    ylabel={Category},
                    xticklabel style={font=\small},
                    yticklabel style={font=\small},
                    scaled y ticks=false,
                    scaled x ticks=false,
                    xbar,
                    ytick={1,2,3},
                    yticklabels={Group 1,Group 2,Group 3},
                    nodes near coords,
                    enlarge x limits={upper, value=0.18},
                    xmin=0,
                    width=.43\textwidth,
                    height=.48\textheight
                ]
                \nextgroupplot[title={\normalsize \textbf{Configuration 1}}]
                    \addplot[draw=black, fill=gray!50, bar width=4.5pt]
                        table[x=C1_Method_A, y expr=\coordindex+1]{\dataComparison};
                    \addplot[draw=black, fill=white, bar width=4.5pt]
                        table[x=C1_Method_B, y expr=\coordindex+1]{\dataComparison};
                    \coordinate (bar top left) at (rel axis cs:0,1);

                \nextgroupplot[
                        title={\normalsize \textbf{Configuration 2}},
                        legend style={
                            at={($(0,0)+(1cm,1cm)$)},
                            legend columns=2,
                            fill=none,
                            draw=none,
                            anchor=south,
                            inner sep=0.2em,
                            column sep=5pt,
                            align=center
                        },
                        legend to name=basicBarLegend
                    ]
                    \addplot[draw=black, fill=gray!50, bar width=4.5pt]
                        table[x=C2_Method_A, y expr=\coordindex+1]{\dataComparison};
                    \addlegendentry{Method A}
                    \addplot[draw=black, fill=white, bar width=4.5pt]
                        table[x=C2_Method_B, y expr=\coordindex+1]{\dataComparison};
                    \addlegendentry{Method B}
                    \coordinate (bar top right) at (rel axis cs:1,1);
            \end{groupplot}

            \coordinate (bar legend middle) at ($(bar top left)!.5!(bar top right)$);
            \coordinate (bar legend position)
                at (bar legend middle |- current bounding box.north);
            \node[above] at (bar legend position)
                {\pgfplotslegendfromname{basicBarLegend}};
        \end{tikzpicture}
    \end{adjustbox}
\end{frame}

% ESEMPIO AVANZATO: personaggi e fumetti riusati dalla tecnica di Nicola.
% Dimostra testo sostituito, due speaker simultanei e una scena finale,
% sincronizzati con l'highlighting del grafo.
% Il grafo incluso a sinistra e' la stessa figura condivisa dell'esempio
% basic: qui riceve intervalli diversi e quindi colorazioni diverse, sincronizzate
% con la comparsa dei personaggi e dei fumetti.

\begin{frame}
    \frametitle{08 Advanced: synchronized characters and speech}

    \newcommand{\stepA}{1}
    \newcommand{\stepBOne}{2}
    \newcommand{\stepBTwo}{3}
    \newcommand{\stepC}{4-5}
    \newcommand{\stepAAndC}{5}

    \begin{columns}[T]
        \begin{column}{.42\textwidth}
            \begin{center}
                \begin{tikzpicture}[node distance=2.5em and 2.8em]
    \node[draw, rounded corners, fill=white,
        on slide={<\stepA,\stepAAndC> highlight light}] (a) {Step A};
    \node[draw, rounded corners, fill=white, below left=of a,
        on slide={<\stepBOne> highlight light}] (b1) {Step B1};
    \node[draw, rounded corners, fill=white, below right=of a,
        on slide={<\stepBTwo> highlight light}] (b2) {Step B2};
    \node[draw, rounded corners, fill=white, below=of b1,
        on slide={<\stepC,\stepAAndC> highlight light}] (c) {Step C};

    \draw[->, link,
        on slide={<\stepA,\stepAAndC> highlight link}]
        (a) -- node[on link, on slide={<\stepA,\stepAAndC> highlight text}] {branch 1} (b1);
    \draw[->, link,
        on slide={<\stepA,\stepAAndC> highlight link}]
        (a) -- node[on link, on slide={<\stepA,\stepAAndC> highlight text}] {branch 2} (b2);
    \draw[->, link,
        on slide={<\stepBOne> highlight link}]
        (b1) -- node[on link, on slide={<\stepBOne> highlight text}] {merge} (c);
    \draw[->, link,
        on slide={<\stepBTwo> highlight link}]
        (b2) -- node[on link, on slide={<\stepBTwo> highlight text}] {merge} (c);
\end{tikzpicture}

            \end{center}
        \end{column}
        \begin{column}{.58\textwidth}
            \begin{center}
                \begin{tikzpicture}
                    \node[dave, minimum size=.8cm, visible on=<1-4>]
                        (dave-first) at (0,1.2) {};
                    \node[speaking question, callout absolute pointer=(dave-first.mouth),
                        anchor=west, visible on=<1>]
                        at ([xshift=1cm]dave-first.east) {\alert{Speaker 1:} Speech 1};
                    \node[speaking question, callout absolute pointer=(dave-first.mouth),
                        anchor=west, visible on=<2>]
                        at ([xshift=1cm]dave-first.east) {\alert{Speaker 1:} Speech 2};
                    \node[speaking question, callout absolute pointer=(dave-first.mouth),
                        anchor=west, visible on=<3>]
                        at ([xshift=1cm]dave-first.east) {\alert{Speaker 1:} Speech 3};
                    \node[speaking question, callout absolute pointer=(dave-first.mouth),
                        anchor=west, visible on=<4>]
                        at ([xshift=1cm]dave-first.east) {\alert{Speaker 1:} Speech 4};

                    \node[judge, minimum size=.8cm, visible on=<3-4>]
                        (judge-second) at (0,-1.2) {};
                    \node[speaking question, callout absolute pointer=(judge-second.mouth),
                        anchor=west, visible on=<3>]
                        at ([xshift=1cm]judge-second.east) {\alert{Speaker 2:} Speech 1};
                    \node[speaking question, callout absolute pointer=(judge-second.mouth),
                        anchor=west, visible on=<4>]
                        at ([xshift=1cm]judge-second.east) {\alert{Speaker 2:} Speech 2};

                    \node[dave, minimum size=.8cm, visible on=<5>]
                        (dave-final) at (0,0) {};
                    \node[speaking question, callout absolute pointer=(dave-final.mouth),
                        anchor=west, visible on=<5>]
                        at ([xshift=1cm]dave-final.east) {\alert{Final speaker:} Final speech};
                \end{tikzpicture}
            \end{center}
        \end{column}
    \end{columns}
\end{frame}

% ESEMPIO AVANZATO: riga di chevron con descrizioni e lista in primo piano.
% Riprende la struttura della tecnica originale di Nicola:
% - stile arrow per i chevron;
% - stile desc per le schede sotto ciascun chevron;
% - dissolvenza delle schede con text opacity e fill opacity;
% - lista in primo piano con visible on e highlighting progressivo.

\begin{frame}
    \frametitle{09 Advanced: chevrons and fading background}

    \pgfmathsetmacro\numberFadeStart{int(2)}
    \pgfmathsetmacro\numberItemA{int(\numberFadeStart+0)}
    \pgfmathsetmacro\numberItemB{int(\numberItemA+1)}
    \pgfmathsetmacro\numberItemC{int(\numberItemB+1)}
    \pgfmathsetmacro\numberFadeEnd{int(\numberItemC+0)}

    \begin{center}
        A sentence introducing the \alert{three phases}\ldots
    \end{center}

    \vspace{1em}

    \begin{adjustbox}{max totalsize={.9\textwidth}{\textheight},center}
        \begin{tikzpicture}
    \begin{scope}

        % Riga di frecce a forma di chevron (stile arrow, forma signal), una accanto all'altra a 3pt di distanza
        % (right=3pt of ...). Nelle figure di Nicola è fatta con la libreria "chains" (start chain, on chain), che
        % qui non è caricata perché rovina altre figure del template.
        \node[arrow] (phase 1) {\textbf{Phase 1}};
        \node[arrow, right=3pt of phase 1] (phase 2) {\textbf{Phase 2}};
        \node[arrow, right=3pt of phase 2] (phase 3) {\textbf{Phase 3}};

        % Sotto ciascuna, una scheda con la descrizione (stile desc, che fissa la larghezza del testo). Dal passo
        % in cui compare l'elenco in primo piano le schede si schiariscono: text opacity e fill opacity sono
        % animati con on slide. Le righe sono punti scritti a mano (\textbullet) per non avere i rientri di itemize.
        \node[below=0 of phase 1, desc, align=left, minimum height=2.4cm,
            on slide={<\numberFadeStart-\numberFadeEnd> text opacity=0.1, fill opacity=0.4}] (desc 1) {
            \textbullet~first feature\\
            \textbullet~second feature
        };
        \node[below=0 of phase 2, desc, align=left, minimum height=2.4cm,
            on slide={<\numberFadeStart-\numberFadeEnd> text opacity=0.1, fill opacity=0.4}] (desc 2) {
            \textbullet~first feature\\
            \textbullet~second feature
        };
        \node[below=0 of phase 3, desc, align=left, minimum height=2.4cm,
            on slide={<\numberFadeStart-\numberFadeEnd> text opacity=0.1, fill opacity=0.4}] (desc 3) {
            \textbullet~first feature\\
            \textbullet~second feature\\
            \textbullet~third feature
        };
    \end{scope}

    % L'elenco in primo piano, sopra le schede sbiadite: compare insieme alla dissolvenza (visible on) e
    % ogni punto si evidenzia al suo passo (\alert<n>).
    \node[anchor=north west, font={\normalsize}, visible on=<\numberFadeStart-\numberFadeEnd>] at (desc 1.north west) {%
        \parbox{.62\textwidth}{%
            \textbullet~\alert<\numberItemA>{\textbf{Item 1}}: \alert<\numberItemA>{first short statement}\\
            \textbullet~\alert<\numberItemB>{\textbf{Item 2}}: \alert<\numberItemB>{second short statement}\\
            \textbullet~\alert<\numberItemC>{\textbf{Item 3}}: \alert<\numberItemC>{third short statement}
        }
    };

\end{tikzpicture}

    \end{adjustbox}
\end{frame}

% ESEMPIO AVANZATO: componente parametrico TikZ con \pic.
% La definizione e le chiamate sono riusate direttamente da Nicola in
% examples/assets/figures/fig_pic.tex. Il componente riceve nome, posizione e overlay:
% la stessa costruzione puo' quindi essere istanziata piu' volte senza
% duplicare il disegno interno.

\begin{frame}
    \frametitle{10 Advanced: reusable component with \texttt{\textbackslash pic}}

    \begin{columns}[T]
        \begin{column}{.38\textwidth}
            The same component is instantiated twice:

            \vspace{.5em}

            \begin{itemize}\setlength{\itemsep}{.7em}
                \item \textbf{Name}: \texttt{a} or \texttt{b}; it becomes part of the node names.
                \item \textbf{Position}: the second argument moves the whole group.
                \item \textbf{Overlay}: the third argument selects when that group is highlighted.
                \item The relation uses the generated nodes
                    \texttt{box-a-3} and \texttt{box-b-3}.
            \end{itemize}
        \end{column}
        \begin{column}{.62\textwidth}
            \begin{adjustbox}{width=\linewidth, max height=.72\textheight, center}
                \begin{tikzpicture}
    % Un componente riusabile della figura: pics/<nome>/.style n args={3}.
    % Argomenti: #1 = nome (serve a dare nomi unici ai nodi), #2 = posizione, #3 = intervalli di overlay in cui
    % il gruppo è evidenziato (anche non contigui). Si chiama con \pic{box group={nome}{(x,y)}{intervalli}};
    \tikzset{
        pics/box group/.style n args={3}{
            code={
                \begin{scope}[shift={#2}]
                    \foreach \i in {1,2,3}{
                        \node[rboxes, minimum width=4em,
                            on slide={<#3> highlight light}] (box-#1-\i) at (0,{-0.9*\i}) {#1.\i};
                    }
                \end{scope}
            }
        }
    }

    \pic{box group={a}{(0,0)}{2-2, 4-4}};
    \pic{box group={b}{(5,0)}{3-4}};

    % I nodi creati dal componente si usano come qualunque altro nodo: box-<nome>-<numero>.
    \draw[->, link, on slide={<4-4> highlight link}] (box-a-3) -- node[on link] {relation} (box-b-3);
\end{tikzpicture}

            \end{adjustbox}
        \end{column}
    \end{columns}
\end{frame}

% ESEMPIO AVANZATO: tre tecniche di callout nella stessa figura.
% La figura e i relativi stili sono riusati direttamente da Nicola in
% examples/assets/figures/fig_callouts.tex:
% 1. callout con codice, un puntatore e linguetta;
% 2. callout con più puntatori verso più nodi;
% 3. personaggio con fumetto.
% I tre casi compaiono su overlay successivi tramite visible on e on slide.

\begin{frame}
    \frametitle{11 Advanced: callouts and speech bubbles}

    \begin{adjustbox}{width=.95\linewidth, max height=.8\textheight, center}
        \begin{tikzpicture}
    % Stili per un callout con PIÙ puntatori (da tex.stackexchange.com, domanda 432276): store path salva il
    % percorso del nodo, callouts={(punto1), (punto2), ...} disegna una "coda" verso ciascun punto.
    \tikzset{
        % senza coda: rettangolo senza riempimento, serve a disegnare il callout con più code (vedi sotto)
        basic question no callout/.style={
            rectangle,
            preaction={draw=none, thick},
            fill=none
        },
        add path/.style={
            decoration={show path construction,
                moveto code={\xdef\savedpath{\savedpath (\tikzinputsegmentfirst)}},
                lineto code={\xdef\savedpath{\savedpath -- (\tikzinputsegmentlast)}},
                curveto code={\xdef\savedpath{\savedpath .. controls (\tikzinputsegmentsupporta) and (\tikzinputsegmentsupportb) ..(\tikzinputsegmentlast)}},
                closepath code={\xdef\savedpath{\savedpath -- cycle}}
            },
            decorate
        },
        store path/.style={add path},
        store path/.prefix code={\xdef\savedpath{}},
        callouts/.style={
            store path,
            append after command={
                foreach \target in {#1}{
                    ($(callout)!2pt!-90:\target$)--\target --($(callout)!2pt!90:\target$)
                } \savedpath
            },
            alias=callout
        }
    }

    \node[rboxes] (a) at (0,0) {Element A};
    \node[rboxes] (b) at (4,0) {Element B};
    \node[rboxes] (c) at (8,0) {Element C};

    % Passo 1: callout con del codice dentro e una linguetta col titolo (stile box label). Un solo puntatore:
    % callout absolute pointer. Il codice è testo a larghezza fissa; con minted (vedi sotto) avrebbe i colori.
    \node[basic question no color, callout absolute pointer={(a.north)}, anchor=south,
        visible on=<1>] (code callout) at ([yshift=1.3cm]a.north east) {%
        \parbox{.32\textwidth}{\ttfamily\scriptsize
            \{\\
            ~~"key": "value",\\
            ~~"status": 1\\
            \}}%
    };
    \node[box label, visible on=<1>] at (code callout.north) {\textbf{Label tab}};

    % Passo 2: un unico callout con tre code, verso tre nodi (stile callouts, definito sopra).
    % Il disegno (bordo e riempimento) è applicato solo al passo 2 con on slide: così anche le code, che fanno
    % parte dello stesso percorso, compaiono solo a quel passo.
    \path[basic question no callout, on slide={<2> preaction={draw, very thick}, fill=VeryLightGrey!50}]
        node[callouts={(a.south), (b.south), (c.south)}, visible on=<2>] at (4,-2.2) {%
            \parbox{.3\textwidth}{One bubble with several pointers}%
        };

    % Passo 3: personaggio (tikzpeople) con bolla di dialogo; l'ancora "mouth" è la bocca del personaggio.
    \node[devil, minimum size=.6cm, visible on=<3>] (attacker) at (10.5,-1.8) {};
    \node[speaking question, callout absolute pointer={(attacker.mouth)}, anchor=east,
        visible on=<3>] at ([xshift=-.6cm, yshift=-.3cm]attacker.west) {%
        \parbox{.2\textwidth}{A speech bubble from a character}%
    };

\end{tikzpicture}

    \end{adjustbox}
\end{frame}

% ESEMPIO AVANZATO: titolo con autori, loghi e finanziamenti.
% Il gruppo locale mantiene i metadati dell'esempio separati da quelli della
% presentazione principale.

\begingroup
\title{12 Advanced: title slide with logos and funding}
\author{\small \textbf{Speaker}\inst{1} \and Author Two\inst{1} \and Author Three\inst{2}}
\institute{\inst{1} Affiliation 1 \and \inst{2} Affiliation 2}
\date{}

\begin{frame}
    \titlepage

    \vfill

    \vspace{-2em}
    \begin{adjustbox}{max totalsize={\textwidth}{.125\textheight},center}
        \begin{tabular}{lr}
            \includegraphics[width=.05\textwidth]{assets/img/patrick_happy.jpg} &
            \includegraphics[width=.1\textwidth]{assets/img/patrick_scientist.jpg}
        \end{tabular}
    \end{adjustbox}

    \parMini{The work was partially supported by
        \emph{i)} project A, funded by the funder under the call number 1;
        \emph{ii)} project B, funded by another funder.
    }
\end{frame}
\endgroup


% Le due slide originali di Nicola hanno dipendenze specifiche, caricate
% solo quando queste slide legacy vengono incluse.
% Legacy Nicola: macro specifiche alla notazione e agli esperimenti originali.
% Gli helper generici sono definiti in 00_macro_common.tex e non vengono duplicati qui.

\newcommand*{\prop}{\ensuremath{p}}
\newcommand*{\target}{\ensuremath{ToC}}
\newcommand*{\targetSet}{\referred{\setStyle{TOC}}}
\newcommand*{\mechanism}{\ensuremath{m}}
% \newcommand*{\evidCollectionModel}{\ensuremath{\mathcal{E}}}
\newcommand*{\evidCollectionModel}{\ensuremath{E}}
\newcommand*{\evalFn}{\ensuremath{\mathcal{F}}}
% \newcommand*{\certModel}{\ensuremath{\mathcal{CM}}}
\newcommand*{\certModel}{\ensuremath{CM}}
\newcommand*{\certModelSet}{\referred{\setStyle{\certModel}}}
%\newcommand*{\cert}{\ensuremath{\mathcal{C}}}
\newcommand*{\cert}{\ensuremath{C}}
\newcommand*{\certSet}{\referred{\setStyle{\cert}}}
\newcommand*{\evid}{\ensuremath{ev}}
%\newcommand*{}{def}
\newcommand*{\attr}{\ensuremath{attr}}
\newcommand*{\repo}{\ensuremath{\mathcal{R}}}

\newcommand*{\signature}{\texttt{sig}}
\newcommand*{\signatureCA}[1][]{\signature\ensuremath{_{\text{CA}_{\text{#1}}}}}
\newcommand*{\signatureAL}[1][]{\signature\ensuremath{_{\text{AL}_{\text{#1}}}}}

\newcommand*{\dimension}{\ensuremath{d}}
\newcommand{\verificationStep}{\ensuremath{v}}

\newcommand{\targetInput}{\ensuremath{\mathcal{I}}}

%\newcommand{\dimension}{\ensuremath{d}}

\newcommand{\abbrvPreTraining}{\text{\textit{pt}}}
\newcommand{\abbrvAdaptation}{\text{\textit{ad}}}
\newcommand{\abbrvInference}{\text{\textit{in}}}
\newcommand{\abbrvArtifact}{\text{\textit{ar}}}

\newcommand{\viewPreTraining}{\ensuremath{\view_{\abbrvPreTraining}}}
\newcommand{\viewAdaptation}{\ensuremath{\view_{\abbrvAdaptation}}}
\newcommand{\viewInference}{\ensuremath{\view_{\abbrvInference}}}
\newcommand{\viewArtifact}{\ensuremath{\view_{\abbrvArtifact}}}
\newcommand{\viewAdaptationInference}{\ensuremath{\view_{\abbrvAdaptation+\abbrvInference}}}
\newcommand{\dimensionPreTraining}{\ensuremath{\dimension_{\abbrvPreTraining}}}
\newcommand{\dimensionAdaptation}{\ensuremath{\dimension_{\abbrvAdaptation}}}
\newcommand{\dimensionInference}{\ensuremath{\dimension_{\abbrvInference}}}
\newcommand{\dimensionAdaptationInference}{\ensuremath{\dimension_{\abbrvAdaptation+\abbrvInference}}}
\newcommand{\dimensionArtifact}{\ensuremath{\dimension_{\abbrvArtifact}}}

\newcommand{\succSymbol}{\ensuremath{\text{\textit{Succ}}}}
\newcommand{\succFn}[1]{\ensuremath{\succSymbol(#1)}}

\newcommand{\foundation}{\text{found}}

\newcommand{\testCaseOutput}{\ensuremath{out}}
\newcommand{\testCaseResult}{\ensuremath{res}}


%\newcommand{\hypergraphEvalRule}{\ensuremath{H_{\text{eval}}}}\
\newcommand{\hypergraphEvalRule}{\ensuremath{\mathcal{R}}}
\newcommand{\hypergraphAll}{\ensuremath{H_{\text{cert}}}}
\newcommand{\hypergraphInstance}{\ensuremath{\mathcal{H}}}
\newcommand{\verificationPath}{\ensuremath{\mathcal{V}}}

\newcommand{\hypergraphCertModel}{\ensuremath{H}}
%\newcommand{\hypergraphCert}{\ensuremath{H_{\text{cert}}}}
\NewDocumentCommand{\hypergraphCert}{o}{%
  \ensuremath{H_{\text{cert}\IfValueT{#1}{, #1}}}%
}

\newcommand{\hyperpath}{\ensuremath{\text{\emph{path}}}}

\newcommand{\FPathStarSymbol}{\ensuremath{\mathcal{F}}}
\newcommand{\BPathStarSymbol}{\ensuremath{\mathcal{B}}}
\newcommand{\BFPathStarSymbol}{\ensuremath{\mathcal{BF}}}
\newcommand{\FPathStarOf}[1]{\ensuremath{\FPathStarSymbol{(#1)}}}
\newcommand{\BPathStarOf}[1]{\ensuremath{\BPathStarSymbol{(#1)}}}
\newcommand{\BFPathStarOf}[1]{\ensuremath{\BFPathStarSymbol{(#1)}}}
\newcommand{\PathStarSymbol}{\ensuremath{\text{\emph{PathsSimple}}}}
\newcommand{\PathStarOf}[1]{\ensuremath{\PathStarSymbol{(#1)}}}

\newcommand{\evalNode}{\ensuremath{\evalFn}}
\newcommand{\evalNodeAll}{\ensuremath{\widehat{\evalFn}}}

% \newcommand{\nodes}{\ensuremath{V}}
% \newcommand{\edges}{\ensuremath{E}}
\newcommand{\nodes}{\ensuremath{N}}
\newcommand{\arcs}{\ensuremath{A}}

%\newcommand{\node}{\ensuremath{n}}

\newcommand{\modelName}[1]{\texttt{#1}}

\newcommand{\llmGemma}[1][short]{%
  \ifthenelse{\equal{#1}{short}}{%
    \modelName{gemma3:27b}%
  }{%
    \modelName{gemma3:27b-it-q4\_K\_M}%
  }%
}
\newcommand{\llmLlama}[1][short]{%
  \ifthenelse{\equal{#1}{short}}{%
    \modelName{llama3.3:70b}%
  }{%
    \modelName{llama3.3:70b-instruct-q5\_K\_M}%
  }%
}
\newcommand{\llmDeepseek}[1][short]{%
  \ifthenelse{\equal{#1}{short}}{%
    \modelName{deepseek-r1:70b}%
  }{%
    \modelName{deepseek-r1:70b}%
  }%
}
\newcommand{\llmMistral}[1][short]{%
  \ifthenelse{\equal{#1}{short}}{%
    \modelName{mistral:123b}%
  }{%
    \modelName{mistral-large:123b-instruct-2411-q5\_K\_M}%
  }%
}
\newcommand{\llmQwenCoder}[1][short]{%
  \ifthenelse{\equal{#1}{short}}{%
    \modelName{qwen3-coder:30b}%
  }{%
    \modelName{qwen3-coder:30b-a3b-q4\_K\_M}%
  }%
}
\newcommand{\llmGptOss}[1][short]{%
  \ifthenelse{\equal{#1}{short}}{%
    \modelName{gpt-oss:20b}%
  }{%
    \modelName{gpt-oss:20b}%
  }%
}
\newcommand{\llmCodeLlama}[1][short]{%
  \ifthenelse{\equal{#1}{short}}{%
    \modelName{codellama:70b}%
  }{%
    \modelName{codellama:70b-instruct-q5\_K\_M}%
  }%
}

% \newcommand{\llmGemma}{\modelName{gemma3:27b-it-q4\_K\_M}}
% \newcommand{\llmLlama}{\modelName{llama3.3:70b-instruct-q5\_K\_M}}
% \newcommand{\llmDeepseek}{\modelName{deepseek-r1:70b}}
% \newcommand{\llmMistral}{\modelName{mistral-large:123b-instruct-2411-q5\_K\_M}}
% \newcommand{\llmQwenCoder}{\modelName{qwen3-coder:30b-a3b-q4\_K\_M}}
% \newcommand{\llmGptOss}{\modelName{gpt-oss:20b}}
% \newcommand{\llmCodeLlama}{\modelName{codellama:70b-instruct-q5\_K\_M}}


\newcommand{\E}{\ensuremath{E}}
\newcommand{\Efirst}{\ensuremath{\E_{\text{first}}}}
\newcommand{\Elast}{\ensuremath{\E_{\text{last}}}}
\newcommand{\Erandom}{\ensuremath{\E_{\text{random}}}}
\newcommand{\Enull}{\ensuremath{\E_{\text{null}}}}

%\newcommand{\Prompts}{\ensuremath{\mathbb{P}}}

\newcommand{\llm}{\ensuremath{\text{\emph{app}}}}

\newcommand{\probe}{\ensuremath{pr}}
\newcommand{\myN}{\ensuremath{n}}
\newcommand{\myM}{\ensuremath{m}}
\newcommand{\accordance}{\ensuremath{\text{\emph{res}}}}
%\newcommand{\probeSet}{\ensuremath{P}}
\newcommand{\prompt}{\ensuremath{\text{\emph{prompt}}}}

\newcommand{\answer}{\ensuremath{\text{\emph{ans}}}}

\newcommand{\orderBias}{\ensuremath{\text{\emph{bias}}}}

\newlength{\oldtabcolsep}
\setlength{\oldtabcolsep}{\tabcolsep}

\newcommand{\tailSetSymbol}{\ensuremath{\text{\emph{tail}}}}
\newcommand{\headSetSymbol}{\ensuremath{\text{\emph{head}}}}

\newcommand{\node}{\ensuremath{\text{\emph{node}}}}
\newcommand{\arc}{\ensuremath{\text{\emph{arc}}}}

\newcommand{\headSet}[1]{\ensuremath{\headSetSymbol(#1)}}
\newcommand{\tailSet}[1]{\ensuremath{\tailSetSymbol(#1)}}

\newcommand{\guardFunctionSymbol}{\ensuremath{G}}
\newcommand{\guardFunction}[1]{\ensuremath{\guardFunctionSymbol(#1)}}
\newcommand{\evidenceFunctionSymbol}{\ensuremath{EV}}
\newcommand{\evidenceFunction}[1]{\ensuremath{\evidenceFunctionSymbol(#1)}}
\newcommand{\nil}{\texttt{nil}}
\newcommand{\error}{\texttt{error}}

%%%%%%%%%
% new macro

\newcommand{\propertyLabel}{\ensuremath{p}}
\newcommand{\behaviorSet}{\ensuremath{\mathcal{B}}}
\newcommand{\verificationStepTarget}{\ensuremath{t}}
\newcommand{\verificationStepActivity}{\ensuremath{s}}
\newcommand{\conditionGuardSymbol}{\ensuremath{\text{\emph{g}}}}
\newcommand{\conditionGuardFunction}[1]{\ensuremath{\conditionGuardSymbol(#1)}}

%\newcommand{\verificationPathAnnotated}{\ensuremath{\overline{\verificationPath}}}

\newcommand{\verificationGraphAnnotated}{\ensuremath{\overline{\behaviorSet}}}
\newcommand{\verificationNodesAnnotated}{\ensuremath{\overline{\nodes}}}
\newcommand{\verificationArcsAnnotated}{\ensuremath{\overline{\arcs}}}


\newcommand{\sourceFunctionSymbol}{\ensuremath{\text{\emph{source}}}}
\newcommand{\sourceFunction}[1]{\ensuremath{\sourceFunctionSymbol(#1)}}

\newcommand{\firstFunctionSymbol}{\ensuremath{\text{\emph{first}}}}
\newcommand{\firstFunction}[1]{\ensuremath{\firstFunctionSymbol(#1)}}

\newcommand{\lastFunctionSymbol}{\ensuremath{\text{\emph{last}}}}
\newcommand{\lastFunction}[1]{\ensuremath{\lastFunctionSymbol(#1)}}

\newcommand{\arcCounterSymbol}{\ensuremath{c}}
\newcommand{\arcCounterOf}[1]{\ensuremath{\arcCounterSymbol[#1]}}



\newcommand{\applTypeTextStandalone}{Type 1: Standalone}
\newcommand{\applTypeTextIntegratedArtifactGeneration}{Type 2: Integrated -- Artifact Generation}
\newcommand{\applTypeTextIntegratedDecisionMaking}{Type 3: Integrated -- Decision-Making}

\newcommand{\probeOutput}{\ensuremath{\text{\emph{out}}}}
\newcommand{\probeOutputTrue}{\ensuremath{\probeOutput^{\top}}}
\newcommand{\probeOutputFalse}{\ensuremath{\probeOutput^{\bot}}}

\newcommand{\FKGLSymbolFunction}{\ensuremath{\text{\emph{FK}}}}
\newcommand{\FKGLFunction}[1]{\ensuremath{\FKGLSymbolFunction(#1)}}
\newcommand{\DLSymbolFunction}{\ensuremath{\text{\emph{DL}}}}
\newcommand{\DLFunction}[1]{\ensuremath{\DLSymbolFunction(#1)}}

\newcommand{\arcsCardinality}{\ensuremath{\vert \arcs \vert}}
\newcommand{\nodesCardinality}{\ensuremath{\vert \nodes \vert}}

\newcommand{\nodesFailed}{\ensuremath{\overline{F}}}


\newcommand{\fullCiteMini}[1]{\parMini{\fullcite{#1}}}


\newcommand{\fullCiteAlmostMini}[1]{\parAlmostMini{\fullcite{#1}}}

\newcommand*{\service}[1]{\ensuremath{s_{#1}}}

\newcommand*{\pavail}{\ensuremath{\prop_\text{avail}}}
\newcommand*{\ppriv}{\ensuremath{\prop_{\text{priv}}}}
\newcommand*{\pConfTransit}{\ensuremath{\prop_{\text{conf\_net}}}}
\newcommand*{\pConfStorage}{\ensuremath{\prop_{\text{conf\_str}}}}


%\newcommand*{\relwordTargetCloudTrad}{def}
\newcommand*{\abbrvServiceStorage}{\text{str}}
\newcommand*{\abbrvServiceDatabase}{\text{db}}
\newcommand*{\abbrvServiceAuth}{\text{auth}}
\newcommand*{\abbrvServiceNetwork}{\text{net}}
\newcommand*{\abbrvServiceKube}{\text{k8s}}
\newcommand*{\abbrvServiceAnalytics}{\text{dawp}}
% \newcommand*{\abbrvServiceNetworkInt}{\text{netinf}}
% \newcommand*{\abbrvServiceStorageInt}{\text{strinf}}

\newcommand*{\abbrvServiceNetworkIntCloud}{\text{netinf\_cl}}
\newcommand*{\abbrvServiceStorageIntCloud}{\text{strinf\_cl}}
\newcommand*{\abbrvServiceNetworkIntEdge}{\text{netinf\_edge}}
\newcommand*{\abbrvServiceStorageIntEdge}{\text{strinf\_edge}}

\newcommand*{\abbrvMechanismUserAuth}{\text{user}}
\newcommand*{\abbrvMechanismScheduling}{\text{sched}}
\newcommand*{\abbrvMechanismData}{\text{data}}

\newcommand*{\abbrvHandlerApp}{\text{app}}
\newcommand*{\abbrvHandlerCSP}{\text{csp}}
\newcommand*{\abbrvHandlerISP}{\text{isp}}

\newcommand*{\CA}{\emph{\gls{CA}}}
\newcommand*{\CAApplication}{\CA\ensuremath{_{\abbrvHandlerApp}}}
\newcommand*{\CACSP}{\CA\ensuremath{_{\abbrvHandlerCSP}}}
\newcommand*{\CAISP}{\CA\ensuremath{_{\abbrvHandlerISP}}}

\newcommand*{\AL}{\emph{\gls{AL}}}
\newcommand*{\ALApplication}{\AL\ensuremath{_{\abbrvHandlerApp}}}
\newcommand*{\ALCSP}{\AL\ensuremath{_{\abbrvHandlerCSP}}}
\newcommand*{\ALISP}{\AL\ensuremath{_{\abbrvHandlerISP}}}

\newcommand{\approachSOTA}{\text{\emph{SOTA}}}
\newcommand{\approachOur}{\text{\emph{multi}}}
\newcommand{\certTimeSymbol}{\ensuremath{T}}
\newcommand{\certTimeFunction}[1]{\ensuremath{\certTimeSymbol}(#1)}

\newcommand*{\certStatusStyle}[1]{\texttt{#1}}
\newcommand*{\certStatusValid}[1]{\certStatusStyle{valid(}#1\certStatusStyle{)}} % chktex 9 chktex 36
\newcommand*{\certStatusSuspended}[1]{\certStatusStyle{suspended(}#1\certStatusStyle{)}} % chktex 9 chktex 36
\newcommand*{\certStatusRevoked}[1]{\certStatusStyle{revoked(}#1\certStatusStyle{)}} % chktex 9 chktex 36

% Legacy Nicola: colori, helper e stili specifici delle figure originali.
% Gli stili generici sono definiti in 00_tikz_common.tex e non vengono duplicati qui.

\colorlet{PoisonedRedForText}{white!70!red!100}
\colorlet{PoisonedRedOut}{red!70}
\colorlet{UnPoisonedGreen}{green!40!black!40}
\colorlet{UnPoisonedGreenForText}{green!30!white!90}
\definecolor{NotSoDarkGrey}{RGB}{160, 160, 160}
\definecolor{Completed}{RGB}{82,163,0}
\definecolor{InProgress}{RGB}{204, 102, 0}
\definecolor{colorBlindBiOrange}{HTML}{E66100}
\definecolor{colorBlindBiViolet}{HTML}{5D3A9B}

\newcommand{\normalwidth}{0.3pt}
\newcommand{\thickwidth}{0.8pt}
\newcommand{\forhash}{for hash}
\newcommand{\forrisk}{for risk}
\newcommand{\forToolTraining}{for tool training}
\newcommand{\forToolTesting}{for tool testing}
\newcommand{\poisonedRF}{1}
\newcommand{\headerEnsemblePoisoned}{Hash\-/Based Ensemble (Poisoned Training Sets)}
\newcommand{\headerEnsembleNonPoisoned}{Hash\-/Based Ensemble (Clean Training Sets)}
\newcommand{\headerEnsembleDefault}{Hash\-/Based Ensemble}
\newcommand{\headerRfTestingPoisoned}{Poisoned}
\newcommand{\headerRfTestingNonPoisoned}{Non\-/Poisoned}
\newcommand{\headerRfOther}{}
\newcommand{\codePoisoned}{Poisoned}
\newcommand{\codeNonPoisoned}{Non\-/Poisoned}
\newcommand{\devilLarge}{Large}
\newcommand{\devilSmall}{Small}

\newcommand{\includegraphicswithwisibility}[3]{
  \begin{tikzpicture}
    \node[visible with opacity on={#1}] at (0,0) {\includegraphics[#2]{#3}};
  \end{tikzpicture}
}

\tikzset{
  dataset/.style={
    draw,
    cylinder,
    shape border rotate=90,
    align=center,
  },
  rf/.style={
    draw,
    rounded corners,
    fill=white,
  },
  test path/.style={dashed},
  connections/.style={very thin},
  partition/.style={
    draw,
    cylinder,
    shape border rotate=90,
    align=center,
    minimum height=.6cm,
    minimum width=.35cm,
    fill=white,
  },
  poisoning link/.style={
    link,
    decorate,
    decoration={
      amplitude=.5pt,
      segment length=4pt,
      zigzag,
    },
    ultra thin,
  },
  data point/.style={
    draw,
    rounded corners,
    font={\tiny},
    fill=white,
  },
  data point moving/.style={
    circle,
    draw,
    fill=black,
    inner sep=1pt,
  },
  poisoned/.style={fill=PoisonedRedOut},
  majority voting/.style={
    isosceles triangle,
    draw,
  },
  highlight path/.style={
    highlight base,
    draw=LightGoldenRod,
  },
  highlight node/.style={
    highlight base,
    fill=HighlightFill,
  },
  highlight double/.style={
    double=LightGoldenRod,
    double distance=1.2pt,
    line width=.3pt,
  },
  arrow path unpoisoned/.append style={
    double=UnPoisonedGreen,
    arrows={-Latex[fill=UnPoisonedGreen]},
  },
  arrow path poisoned/.append style={
    double=PoisonedRedOut,
    arrows={-Latex[fill=PoisonedRedOut]},
  },
  highlight path poisoned/.style={},
  highlight path unpoisoned/.style={},
  invisible with opacity/.style={opacity=0.025},
  visible with opacity on/.style={alt=#1{}{invisible with opacity}},
  highlight line in pgfplot/.style={
    LightGoldenRod,
    line width=3pt,
  },
  poisoning attack/.style={
    draw,
    fill=white,
  },
  super container/.style={
    draw,
    dashed,
  },
  low link/.style={
    link,
    ultra thin,
  },
  bar/.style={
    thin,
    dotted,
  },
  from test/.style={ultra thin},
  highlight accuracy/.style={ultra thick},
  explanation/.style={font={\Large}},
  explanation link/.style={>=Latex},
  line am/.style={},
  line m2/.style={dashed},
  service in certification process/.style={
    chamfered rectangle,
    chamfered rectangle angle=30,
    draw,
  },
  connector assurance/.style={
    >=Stealth,
    dashed,
  },
  probe link/.style={
    link,
    dashdotted,
  },
  ca cert/.style={
    entity,
    inner sep=0pt,
  },
  step separator/.style={
    draw,
    thin,
    dotted,
  },
  container line/.style={
    very thin,
    rounded corners,
  },
  label desc/.style={},
  label desc with shape/.style={
    label desc,
    rounded corners,
    draw,
  },
  lattice layer label/.style={
    text width=2.9cm,
    font={\footnotesize},
  },
  not chosen color/.style={black!50},
  lattice layer separator/.style={dashed},
  service label/.style={
    fill=white,
    circle,
    draw,
    very thick,
    inner sep=2.5pt,
  },
  lattice element bound/.style={
    fill=black!15,
    double,
  },
  lattice element selected/.style={
    fill=Goldenrod!15,
    double,
  },
  lattice element not chosen/.style={
    not chosen color,
    densely dotted,
  },
  lattice element not chosen link/.style={
    not chosen color,
    densely dotted,
  },
  service element not chosen/.style={
    text=black!50,
    draw=black!50,
    dashed,
  },
  aggregation block/.style={
    circle,
    draw,
    inner sep=1pt,
  },
  link comment/.style={dashed},
  basic graph node/.style={
    circle,
    draw,
    minimum width=.6cm,
  },
  verification step/.style={
    basic graph node,
    fill=white,
  },
  evaluation function/.style={
    basic graph node,
    double,
    fill=white,
  },
  evaluation function global/.style={
    draw,
    double,
    rectangle,
    minimum width=.6cm,
    minimum height=.6cm,
  },
  hyperarc/.style={>=Latex},
  pics/hyperarc between fan in/.style n args={5}{
    code={
      \path let \p1 = ($(#1)!.5!(#2)$), \p2 = ($(#2)!.5!(#3)$) in
        coordinate (#4) at (\x1, \y2);
      \foreach \src in {#1, #2} {
        \draw[-, hyperarc, #5] (\src) -- (#4);
      }
      \draw[->, hyperarc, #5] (#4) -- (#3);
    }
  },
  pics/hyperarc between fan out/.style n args={5}{
    code={
      \path let \p1 = ($(#1)!.5!(#2)$), \p2 = ($(#2)!.5!(#3)$) in
        coordinate (#4) at (\x1, \y2);
      \foreach \src in {#1, #2} {
        \draw[<-, hyperarc, #5] (\src) -- (#4);
      }
      \draw[-, hyperarc, #5] (#4) -- (#3);
    }
  },
  visited node/.style={},
  visited arc/.style={thick},
  pruned node/.style={fill=gray!30},
  pruned arc/.style={draw=gray!50},
  evidence content/.style={font=\scriptsize},
  orange in violet out double/.style={
    fill=colorBlindBiOrange,
    double=colorBlindBiViolet,
    draw=colorBlindBiViolet,
    very thick,
  },
  orange out/.style={
    draw=colorBlindBiOrange,
    very thick,
  },
  highlighted hyperarc in path/.style={thick},
  hyperarc in verification path/.style={
    very thick,
    highlight path,
  },
  node in verification path/.style={
    very thick,
    highlight path,
  },
  certificate callout/.style={
    rectangle callout,
    fill=VeryLightGrey!50,
    draw,
  },
}

% !TEX root = ../main.tex
% ESEMPIO 0a: slide Background con elenco evidenziato passo passo accanto a una figura.
%   Ogni punto dell'elenco usa \alert<n>, che lo evidenzia solo alla sovrapposizione (overlay) n. Gli stessi
%   numeri sono usati dalla figura (assets/tex/01_fig_cert_process.tex) per evidenziare la parte corrispondente.
%   Da copiare: le macro dei numeri, le due colonne e l'\input della figura dentro adjustbox.

\begin{frame}
    \frametitle{13 Legacy Nicola: highlighted list next to a figure}
    % GUIDA: il concetto, lo schema o il contesto di cui parla il lavoro, con i suoi elementi. Meglio una figura che una definizione
    % lunga: ogni elemento si evidenzia (\alert<n>) mentre lo nomini.

    % Numero di sovrapposizione (overlay) in cui compare ogni passo. Il primo è 2 perché alla
    % sovrapposizione 1 si vede la figura senza evidenziazioni; ogni passo parte dal precedente.
    % Gli stessi nomi sono usati da assets/tex/01_fig_cert_process.tex per evidenziare la figura.
    \pgfmathsetmacro\numberProperty{int(2)}
    \pgfmathsetmacro\numberTarget{int(\numberProperty+1)}
    \pgfmathsetmacro\numberEvidenceCollectionModel{int(\numberTarget+1)}
    \pgfmathsetmacro\numberCertificationModel{int(\numberEvidenceCollectionModel+1)}
    \pgfmathsetmacro\numberExecution{int(\numberCertificationModel+1)}
    \pgfmathsetmacro\numberCertificate{int(\numberExecution+1)}

    \begin{columns}[T]
        \begin{column}{.45\textwidth}

            CHANGE: one sentence introducing the \alert{concept} and its \alert{key words}, followed by the elements

            \vspace{.5em}

            % Un punto per ogni passo: \alert<n> lo evidenzia solo alla sovrapposizione n.
            \begin{itemize}\setlength\itemsep{1em}
                \item \alert<\numberProperty>{CHANGE: step 1}
                \item \alert<\numberTarget>{CHANGE: step 2}
                \item \alert<\numberEvidenceCollectionModel>{CHANGE: step 3}
                \item \alert<\numberCertificationModel>{CHANGE: step 4}
                \item \alert<\numberExecution>{CHANGE: step 5}
                \item \alert<\numberCertificate>{CHANGE: step 6}
            \end{itemize}

        \end{column}
        \begin{column}{.55\textwidth}
            \begin{adjustbox}{max totalsize={\linewidth}{.8\textheight},center}
                
\begin{tikzpicture}

    % This legacy figure predates the template-wide centered-node default.
    \tikzset{every node/.style={font=\scriptsize}}

    \node[entity, on slide={<\numberCertificationModel> highlight light}] (ca) at (-1.5, 0) {CA};

  \node[artifact, on slide={<\numberCertificationModel,\numberExecution> highlight light}] (cm) at (-1.5, -1.65) {\makecell{Certification\\Model}};

    \draw[->, link, on slide={<\numberCertificationModel> highlight link}] (ca) to node[on link, on slide={<\numberCertificationModel> highlight text}]{prepares} (cm);

    \node[artifact, on slide={<\numberProperty-\numberExecution> highlight light}] (properties) at (2.5, -1.65) (property)  {\makecell{Property}};

    \node[service, on slide={<\numberTarget-\numberExecution> highlight light}] (target) at (2.5, -3.25) {\makecell{Target}};

    \draw[->, link, on slide={<\numberCertificationModel-\numberExecution> highlight link}] (cm) to node[on link, on slide={<\numberCertificationModel-\numberExecution> highlight text}] {evaluates} (property);

    %\draw[->, link, on slide={<4,5> highlight link}] (cm) to node[on link, on slide={<4,5> highlight text}] {to assess} (target);
    \draw[->, link, on slide={<\numberTarget-\numberExecution> highlight link}] (property) to node[on link, on slide={<\numberTarget-\numberExecution> highlight text}] {on} (target);

    \node[artifact, on slide={<\numberEvidenceCollectionModel-\numberExecution> highlight light}] (evidence collection model) at (cm |- target) {\makecell{Evidence\\Coll. Model}};

    %\node[entity, inner sep=1pt, on slide={<5> highlight light}] (lab) at (-1.5, -3.65) {\makecell{Accredited\\Lab}};

    \node[entity, inner sep=1pt, on slide={<\numberExecution> highlight light}] (lab) at ([yshift=-45]evidence collection model.south) {\makecell{Accredited\\Lab}};

    \draw[->, link, on slide={<\numberEvidenceCollectionModel, \numberCertificationModel, \numberExecution> highlight link}] (target) to node[on link, on slide={<\numberEvidenceCollectionModel, \numberCertificationModel, \numberExecution> highlight text}] {according to} (evidence collection model);

    %\draw[->, link, on slide={<5> highlight link}] (cm) to node[on link, on slide={<5> highlight text}]{executed by} (lab);
    \draw[->, link, on slide={<\numberExecution> highlight link}] (evidence collection model) to node[on link, on slide={<\numberExecution> highlight text}]{executed by} (lab);

    \draw[->, link, on slide={<\numberExecution> highlight link}] (lab) to node[on link, on slide={<\numberExecution> highlight text}] (against) {against} (target);

    \node[artifact, on slide={<\numberExecution> highlight light}] at ([yshift=-35]against.south) (evidence) {Evidence};

    \draw[->, link, on slide={<\numberExecution> highlight link}] (against) to node[on link, on slide={<\numberExecution> highlight text}] {collecting} (evidence);

    \node[artifact, on slide={<\numberCertificate> highlight light}] at ([yshift=-35]evidence.south) (certificate) {\makecell{Certificate}};

    \draw[ link, ->, on slide={<\numberCertificate> highlight link}] (evidence) to node[on link, on slide={<\numberCertificate> highlight text}]{on success} (certificate);

   %  \draw[link, ->, bend right=70, on slide={<6> highlight link}] (certificate.east) to node[xshift=20, near end, on link, on slide={<6> highlight text}] (includes) {includes} (evidence.east);

   %  \draw[->, link, bend right=45, on slide={<6> highlight link}] (includes) to (target.east);

   %  \draw[->, link, bend right=55, on slide={<6> highlight link}] (includes.east) to (property.east);


\end{tikzpicture}
            \end{adjustbox}
        \end{column}
    \end{columns}

\end{frame}

% !TEX root = ../main.tex
% ESEMPIO 0b: slide Motivations con una figura animata (personaggi e bolle di dialogo).
%   La figura può essere un'immagine qualsiasi. Qui i numeri di sovrapposizione (overlay) sono definiti nel frame:
%   ogni passo parte dal precedente (+1) o resta uguale (+0) quando due cose compaiono insieme.
%   Gli stessi nomi sono usati da assets/tex/02_fig_cert_process.tex.
%   Da copiare: la struttura (macro dei numeri nel frame, figura in assets/tex/, \input dentro adjustbox).
%   La stessa tecnica si riusa nell'approccio, per mostrare una figura con i colori che evidenziano i passi.

\begin{frame}
    \frametitle{14 Legacy Nicola: animated scenario figure}
    % GUIDA: il problema, e perché lo stato dell'arte non basta. Meglio mostrarlo su un esempio concreto (una figura di uno scenario),
    % che può tornare più avanti nell'approccio.

    \pgfmathsetmacro\numberStart{int(1)}
    % passo: si evidenzia l'elemento principale e le sue dipendenze
    \pgfmathsetmacro\numberServiceTopStart{int(\numberStart+0)}
    \pgfmathsetmacro\numberHighlightDependencyStart{int(\numberServiceTopStart+0)}
    % passo: nuova bolla sull'elemento principale, che poi smette di essere evidenziato
    \pgfmathsetmacro\numberServiceTopEvidStart{int(\numberHighlightDependencyStart+1)}
    \pgfmathsetmacro\numberServiceTopEvidEnd{int(\numberServiceTopEvidStart+0)}
    \pgfmathsetmacro\numberServiceTopEnd{int(\numberServiceTopEvidEnd+0)}
    % passo: restano evidenziati gli altri elementi, nuova bolla
    \pgfmathsetmacro\numberCalloutFocusOnOtherStart{int(\numberServiceTopEvidEnd+1)}
    \pgfmathsetmacro\numberCalloutFocusOnOtherEnd{int(\numberCalloutFocusOnOtherStart+0)}
    % passo: doppio evidenziamento di più elementi
    \pgfmathsetmacro\numberHighlightDependencyEnd{int(\numberCalloutFocusOnOtherEnd+0)}
    \pgfmathsetmacro\numberHighlightDependencySomeOkStart{int(\numberCalloutFocusOnOtherStart+0)}
    \pgfmathsetmacro\numberHighlightDependencySomeOkEnd{int(\numberHighlightDependencySomeOkStart+0)}
    % passo: doppio evidenziamento di un altro elemento, nuova bolla
    \pgfmathsetmacro\numberHighlightDependencySomeNotOkStart{int(\numberHighlightDependencySomeOkEnd+1)}
    \pgfmathsetmacro\numberHighlightDependencySomeNotOkStartPartOne{int(\numberHighlightDependencySomeNotOkStart+0)}
    \pgfmathsetmacro\numberHighlightDependencySomeNotOkStartPartOneEnd{int(\numberHighlightDependencySomeNotOkStartPartOne+0)}
    \pgfmathsetmacro\numberHighlightDependencySomeNotOkStartPartTwoStart{int(\numberHighlightDependencySomeNotOkStartPartOneEnd+0)}
    \pgfmathsetmacro\numberHighlightDependencySomeNotOkStartPartTwoEnd{int(\numberHighlightDependencySomeNotOkStartPartTwoStart+0)}
    \pgfmathsetmacro\numberHighlightDependencySomeNotOkEnd{int(\numberHighlightDependencySomeNotOkStartPartTwoEnd+0)}
    \pgfmathsetmacro\numberEnd{int(\numberHighlightDependencySomeNotOkEnd+0)}
    % passo finale: ogni elemento mostra la sua bolla e compare l'immagine
    \pgfmathsetmacro\numberCertStart{int(\numberHighlightDependencySomeNotOkEnd+1)}

    \pgfmathsetmacro\numberPatrickHappy{int(\numberCertStart+1)}
    \pgfmathsetmacro\numberCertEnd{int(\numberPatrickHappy+0)}

    \begin{adjustbox}{max totalsize={.9\linewidth}{\textheight},center}
        
	
\begin{tikzpicture}
			% \tikzset{
			% 	every node/.style={font=\scriptsize}
			% }
	\node[service in certification process,
    on slide={<\numberServiceTopStart> highlight light},
    on slide={<\numberServiceTopEvidStart-\numberServiceTopEvidEnd> highlight light},
    ] (target) {\makecell{Data Analysis Engine (\target)}};
	
    \node[service,below=0.6 of target,
    on slide={<\numberServiceTopStart> highlight light},
    on slide={<\numberCalloutFocusOnOtherStart-\numberCalloutFocusOnOtherEnd> highlight light}
    ] (auth) {Authentication};

	\node[service,below=0.5 of auth,
    on slide={<\numberServiceTopStart> highlight light},
    on slide={<\numberCalloutFocusOnOtherStart-\numberCalloutFocusOnOtherEnd> highlight light}
    ] (db) {Database};

	\node[service,below=0.5 of db,
    on slide={<\numberServiceTopStart> highlight light},
    on slide={<\numberCalloutFocusOnOtherStart-\numberCalloutFocusOnOtherEnd> highlight light}
    ] (ds) {\makecell{Distributed\\Storage}};
	
    \node[service,left=1.5 of ds,
    on slide={<\numberServiceTopStart> highlight light},
    on slide={<\numberCalloutFocusOnOtherStart-\numberCalloutFocusOnOtherEnd> highlight light}
    ] (mn) {\makecell{Mesh\\Network}};
	
    \node[service,right=2 of ds,
    on slide={<\numberServiceTopStart> highlight light},
    on slide={<\numberHighlightDependencySomeNotOkStart-\numberHighlightDependencySomeNotOkEnd> highlight light}
    ] (k8s) {Kubernetes};
	
    \node[service,below=0.7 of mn,
    on slide={<\numberServiceTopStart> highlight light},
    on slide={<\numberCalloutFocusOnOtherStart-\numberCalloutFocusOnOtherEnd> highlight light}
    ] (ni) {\makecell{Network Interface}};
	
    \node[service,below=0.7 of ds,
    on slide={<\numberServiceTopStart> highlight light},
    on slide={<\numberCalloutFocusOnOtherStart-\numberCalloutFocusOnOtherEnd> highlight light}
    ] (si) {\makecell{Storage Interface}};

	\begin{scope}[every node/.style={font=\footnotesize}]
		\draw[->, link] (target) -- node [left] {\tiny depends on} (auth);
		\draw[->, link] ([yshift=-5pt]target.east) -- ++(.25,0) |- node [left,pos=.175] {\tiny depends on} ([yshift=3pt]db.east);
		\draw[->, link] (auth) -- node [left] {\tiny depends on} (db);
		\draw[->, link] (db) -- node [left] {\tiny depends on} (ds);
		\draw[->, link] (ds) -- node [above] {\tiny depends on} (mn);
		\draw[->, link] (target) -| node [above,pos=.25] {\tiny depends on} ([xshift=-5pt]mn.north);
		\draw[->, link] (auth) -| node [above,pos=.25] {\tiny depends on} (mn.north);
		\draw[->, link] (db) -| node [above,pos=.25] {\tiny depends on} ([xshift=5pt]mn.north);
		\draw[->, link] ([yshift=-3pt]k8s.west) -- node [below] {\tiny coordinates} ([yshift=-3pt]ds.east);
		\draw[->, link] ([xshift=-5]k8s.north) |- node [below left] {\tiny coordinates} ([yshift=-3pt]db.east);
		\draw[->, link] (k8s.north) |- node [below left] {\tiny coordinates} (auth.east);
		\draw[->, link] ([xshift=5]k8s.north) |- node [below left] {\tiny coordinates} ([yshift=5pt]target.east);
		\draw[->, link] ([xshift=3]k8s.south) -- ++(0,-.6) -| node [below,pos=.25] {\tiny coordinates} ([xshift=3pt]si.north);
		\draw[->, link] ([xshift=-3]k8s.south) -- ++(0,-.4) -| node [below right] {\tiny coordinates} ([xshift=3pt]ni.north);
		\draw[->, link] (target.east) -- ++(.5,0) |- node [right,pos=.128] {\tiny depends on} ([yshift=3pt]ds.east);
		\draw[->, link] ([xshift=-3]ds.south) -- node [left] {\tiny depends on} ([xshift=-3pt]si.north);
		\draw[->, link] ([xshift=-3]mn.south) -- node [left] {\tiny depends on} ([xshift=-3pt]ni.north);
	\end{scope}


    \path let \p1 = (target), \p2 = (k8s) in
        coordinate (person-1-pos) at (\x2, \y1);
    \coordinate (person-1-pos) at ([xshift=20]person-1-pos); 

    \node[dave, minimum size=.75cm, xshift=20,
        visible on=<\numberStart-\numberEnd>
    ] (person-1) at (person-1-pos) {};
    %
    %
    % 2. let's certify it.
    \node[speaking question, callout absolute pointer=(person-1.mouth), anchor=west,
    visible on=<\numberServiceTopStart>
    ] (person-1-speak-1) at ([xshift=15]person-1.east) {%
        \begin{varwidth}{\widthof{These are its dependencies, }}
            \raggedright
            Let's certify this service. These are its dependencies, specified in \certModel.
        \end{varwidth}%
    };
    %
    %
    % 3. collect evidence from there
    \node[speaking question, callout absolute pointer=(person-1.mouth), anchor=west,
    visible on=<\numberServiceTopEvidStart-\numberServiceTopEvidEnd>
    ] (person-1-speak-2) at ([xshift=15]person-1.east) {%
        \begin{varwidth}{\widthof{These are its dependencies, }}
            \raggedright
            Let's certify this service. First, we collect evidence from its mechanisms.
        \end{varwidth}%
    };
    %
    %
    % 4. exists cert compatible?
    \node[speaking question, callout absolute pointer=(person-1.mouth), anchor=west,
    visible on=<\numberCalloutFocusOnOtherStart-\numberCalloutFocusOnOtherEnd>
    ] (person-1-speak-3) at ([xshift=15]person-1.east) {%
        \begin{varwidth}{\widthof{These are its dependencies, }}
            \raggedright
            Let's certify this service. \alert{These are its dependencies: are they already certified according to $\certModel.\referred{\certModel_i}$}?
        \end{varwidth}%
    };
    %
    %
    % 5. highlight many that are
    % so we need a judge
    \node[judge, minimum size=.75cm, xshift=20,
        %visible on=<\numberHighlightDependencySomeOkStart-\numberHighlightDependencySomeOkEnd>
        visible on=<\numberCalloutFocusOnOtherStart-\numberCalloutFocusOnOtherEnd>
    ] (person-2) at ([yshift=-80]person-1-pos) {};
    \node[speaking question, callout absolute pointer=(person-2.mouth), anchor=west,
    %visible on=<\numberHighlightDependencySomeOkStart-\numberHighlightDependencySomeOkEnd>
        visible on=<\numberCalloutFocusOnOtherStart-\numberCalloutFocusOnOtherEnd>
    ] (person-2-speak-4) at ([xshift=15]person-2.east) {%
        \begin{varwidth}{\widthof{I hereby guarantee that }}
            \raggedright
            I hereby guarantee that these dependencies are already certified.
        \end{varwidth}%
    };
    %
    %
    % 6. pt1 highlight some that are not.
    % we need a judge
    \node[judge, minimum size=.75cm, xshift=20,
    visible on=<\numberHighlightDependencySomeNotOkStart-\numberHighlightDependencySomeNotOkEnd>
    ] (person-5) at (person-1-pos) {};
    \node[speaking question, callout absolute pointer=(person-1.mouth), anchor=west,
    visible on=<\numberHighlightDependencySomeNotOkStart-\numberHighlightDependencySomeNotOkEnd>
    ] (person-2-speak-5) at ([xshift=15]person-1.east) {%
        \begin{varwidth}{\widthof{I hereby declare that}}
            \raggedright
            I hereby declare that these dependencies are not already certified.
        \end{varwidth}%
    };
    % we need dave right below
    \node[dave, minimum size=.75cm, xshift=20,
    visible on=<\numberHighlightDependencySomeNotOkStart-\numberHighlightDependencySomeNotOkEnd>
    ] (person-3) at ([yshift=-80]person-1-pos) {};
    %
    \node[speaking question, callout absolute pointer=(person-3.mouth), anchor=west,
    visible on=<\numberHighlightDependencySomeNotOkStart-\numberHighlightDependencySomeNotOkEnd>
    ] (person-3-speak-5) at ([xshift=15]person-3.east) {%
        \begin{varwidth}{\widthof{Let's repeat}}
            \raggedright
            Let's repeat the process recursively.
        \end{varwidth}%
    };
    %
    %
    % 7. certificate
    % we re-add the judge just because
    \node[judge, minimum size=.75cm, xshift=20,
        visible on=<\numberCertStart-\numberCertEnd>
    ] (person-4) at ([yshift=-50]person-1-pos) {};
    %
    \node[speaking question, callout absolute pointer=(person-4.mouth), anchor=west,
    visible on=<\numberCertStart-\numberCertEnd>
    ] (person-3-speak-5) at ([xshift=15]person-4.east) {%
        \begin{varwidth}{\widthof{Based on the evidence and the }}
            \raggedright
            Based on the evidence and the referred certificates provided to me, I hereby award a \alert{certificate \cert}, including the evidence, the \alert{referred certificates $\{\referred{\cert}\}$}, and my signature. 
        \end{varwidth}%
    };


    % south
    \foreach \src/\pointer in {ni/ni.south, si/si.south}{
        \node[speaking question, callout absolute pointer=(\pointer), anchor=north,
        visible on=<\numberCertStart-\numberCertEnd>,
        %visible on=<\numberCalloutFocusOnOtherStart-\numberCalloutFocusOnOtherEnd>
        ] at ([yshift=-10]\pointer) (\src-cert) {%
            %\begin{varwidth}{\widthof{certified}}
                \tiny I am certified
            %\end{varwidth}
        };
        \node[speaking question, callout absolute pointer=(\pointer), anchor=north,
        %visible on=<\numberCertStart-\numberCertEnd>,
        visible on=<\numberCalloutFocusOnOtherStart-\numberCalloutFocusOnOtherEnd>
        ] at ([yshift=-10]\pointer) (\src-cert-2) {%
            %\begin{varwidth}{\widthof{certified}}
                \tiny I am certified
            %\end{varwidth}
        };
    }
    % mesh
    \node[speaking question, callout absolute pointer=(mn.west), anchor=east,
        visible on=<\numberCertStart-\numberCertEnd>,
        %visible on=<\numberCalloutFocusOnOtherStart-\numberCalloutFocusOnOtherEnd>
        ] (mn-cert) at ([xshift=-10]mn.west) {%
            %\begin{varwidth}{\widthof{certified}}
                \tiny I am certified
            %\end{varwidth}
        };
    \node[speaking question, callout absolute pointer=(mn.west), anchor=east,
        %visible on=<\numberCertStart-\numberCertEnd>,
        visible on=<\numberCalloutFocusOnOtherStart-\numberCalloutFocusOnOtherEnd>
        ] (mn-cert) at ([xshift=-10]mn.west) {%
            %\begin{varwidth}{\widthof{certified}}
                \tiny I am certified
            %\end{varwidth}
        };
    % db
    \path let \p1 = (mn-cert), \p2 = (db) in
        node[speaking question, callout absolute pointer=(db.west), %anchor=east,
         yshift=5,
        visible on=<\numberCertStart-\numberCertEnd>,
        %visible on=<\numberCalloutFocusOnOtherStart-\numberCalloutFocusOnOtherEnd>
        ] (db-cert) at (\x1, \y2) {%
            \tiny I am certified
        };
    % repeat.
    \path let \p1 = (mn-cert), \p2 = (db) in
        node[speaking question, callout absolute pointer=(db.west), %anchor=east,
         yshift=5,
        %visible on=<\numberCertStart-\numberCertEnd>,
        visible on=<\numberCalloutFocusOnOtherStart-\numberCalloutFocusOnOtherEnd>
        ] (db-cert-2) at (\x1, \y2) {%
            \tiny I am certified
        };
    % auth
    \path let \p1 = (db-cert), \p2 = (auth) in
        node[speaking question, callout absolute pointer=(auth.west), %anchor=east,
         yshift=5,
        %visible on=<\numberCalloutFocusOnOtherStart-\numberCalloutFocusOnOtherEnd>
        visible on=<\numberCertStart-\numberCertEnd>
        ] at (\x1, \y2) (auth-cert) {%
            \tiny I am certified
        };
    \path let \p1 = (db-cert), \p2 = (auth) in
        node[speaking question, callout absolute pointer=(auth.west), %anchor=east,
         yshift=5,
        visible on=<\numberCalloutFocusOnOtherStart-\numberCalloutFocusOnOtherEnd>
        %visible on=<\numberCertStart-\numberCertEnd>
        ] at (\x1, \y2) (auth-cert-2) {%
            \tiny I am certified
        };
    % target
    \node[speaking question, callout absolute pointer=(target.north), anchor=south,
        visible on=<\numberCertStart-\numberCertEnd>
        ] at ([yshift=10]target.north) (target-cert) {%
            \tiny I am certified
        };
    % ds
    \node[speaking question, callout absolute pointer=(ds.south east), anchor=north west,
        visible on=<\numberCertStart-\numberCertEnd>,
        %visible on=<\numberCalloutFocusOnOtherStart-\numberCalloutFocusOnOtherEnd>
        ] at (si -| k8s) (ds-cert) {%
            \tiny I am certified
        };
    \node[speaking question, callout absolute pointer=(ds.south east), anchor=north west,
        %visible on=<\numberCertStart-\numberCertEnd>,
        visible on=<\numberCalloutFocusOnOtherStart-\numberCalloutFocusOnOtherEnd>
        ] at (si -| k8s) (ds-cert-2) {%
            \tiny I am certified
        };  
    % k8s
    \node[speaking question, callout absolute pointer=(k8s.east), anchor=west,
        visible on=<\numberCertStart-\numberCertEnd>
        ] at ([xshift=10]k8s.east) (k8s-cert) {%
            \tiny I am certified
        }; 


    \node[visible on=<\numberPatrickHappy>, anchor=north] at ([yshift=-10]person-3-speak-5.south) {\includegraphics[width=.3\textwidth]{assets/img/patrick_happy.jpg}};


\end{tikzpicture}


    \end{adjustbox}

\end{frame}



\end{document}
